% SPDX-License-Identifier: MIT-0
\documentclass[11pt]{article}
\usepackage[T1]{fontenc}
\usepackage[utf8]{inputenc}
\usepackage{lmodern}
\usepackage[a4paper,margin=26mm]{geometry}
\usepackage{amsmath,amssymb,amsthm,mathtools}
\usepackage{microtype,booktabs,array,longtable}
\usepackage[dvipsnames]{xcolor}
\usepackage{enumitem}
\usepackage{fancyhdr}
\usepackage[colorlinks=true,linkcolor=MidnightBlue,citecolor=MidnightBlue,
  urlcolor=MidnightBlue]{hyperref}
\usepackage{bookmark}
\usepackage{listings}
\hypersetup{pdftitle={Multiple-chain preorder polytopes: two proofs of the exponential formula},
 pdfauthor={Research draft prepared for Vladimir Reshetnikov},
 pdfsubject={Athanasiadis--Chapoton Section 9.5; lattice-point enumeration, excursions and slack meanders}}
\setlength{\headheight}{14pt}
\pagestyle{fancy}
\fancyhf{}
\fancyhead[L]{\small Multiple-chain preorder polytopes}
\fancyhead[R]{\small Research draft}
\fancyfoot[C]{\thepage}
\renewcommand{\headrulewidth}{0.3pt}
\setlength{\parskip}{0.25em}
\emergencystretch=2em
\numberwithin{equation}{section}
\newtheorem{theorem}{Theorem}[section]
\newtheorem{lemma}[theorem]{Lemma}
\newtheorem{proposition}[theorem]{Proposition}
\newtheorem{corollary}[theorem]{Corollary}
\theoremstyle{definition}
\newtheorem{definition}[theorem]{Definition}
\newtheorem{example}[theorem]{Example}
\newtheorem{remark}[theorem]{Remark}
\newcommand{\N}{\mathbb N}
\newcommand{\Z}{\mathbb Z}
\newcommand{\Q}{\mathbb Q}
\newcommand{\R}{\mathbb R}
\newcommand{\C}{\mathbb C}
\newcommand{\CT}{\operatorname{CT}_u}
\newcommand{\CTv}{\operatorname{CT}_v}
\newcommand{\supp}{\operatorname{supp}}
\newcommand{\Var}{\operatorname{Var}}
\newcommand{\E}{\mathbb E}
\newcommand{\A}{A_{n;d,c}}
\newcommand{\Adc}[1]{A_{#1;d,c}}
\newcommand{\B}{\mathcal B}
\newcommand{\F}{\mathcal F}
\newcommand{\pt}{\operatorname{pt}}
\newcommand{\pos}[1]{\left[#1\right]_{\geq0}}
\newcommand{\negp}[1]{\left[#1\right]_{<0}}
\newcommand{\abs}[1]{\left|#1\right|}
\newcommand{\Cat}{\operatorname{Cat}}
\DeclareMathOperator{\Ehr}{Ehr}
\lstset{basicstyle=\small\ttfamily,columns=fullflexible,breaklines=true,
 frame=single,rulecolor=\color{gray!45},xleftmargin=0.4em,xrightmargin=0.4em}
\title{\textbf{Multiple-chain preorder polytopes:\\
two proofs of the exponential formula}\\[0.6em]
\large A rectangular extension, unequal block sizes, algebraic generating
functions, fixed-support polynomiality, and support asymptotics}
\author{Research draft prepared for Vladimir Reshetnikov}
\date{20 September 2026}
\begin{document}
\maketitle
\thispagestyle{empty}

\begin{abstract}
Let $H_{n,d}(t)$ enumerate the nonnegative integer points satisfying
$\sum_{i\leq r}\sum_{j\leq d}a_{ij}\leq dr$ for $1\leq r\leq n$, with
$t$ marking nonzero coordinates. We give two complete and structurally
different formal-power-series proofs of
\[
 \sum_{n\geq0}H_{n,d}(t)z^n
 =\exp\left(\sum_{m\geq1}\frac{z^m}{m}
       \sum_{j=0}^{dm}\binom{dm}{j}^{\!2}t^j\right).
\]
This is the conjectural multiple-chain identity of Athanasiadis and Chapoton,
Section~9.5 of arXiv:2605.26916v1, with its normalization removed. The first
proof encodes lattice points by two subsets and applies a constant-term
excursion factorization with a bounded Laurent-polynomial symbol; the second
runs the original coordinate slack process, whose symbol is an unbounded
Laurent series, and evaluates the positive part by stars and bars. Replacing
the budget $dr$ by $cr$ gives the stronger bridge polynomial
$\sum_j\binom{dm}{j}\binom{cm}{j}t^j$; allowing unequal block sizes gives a
commutative formula summed over orders, which is sharp because the support
polynomial is not a function of the multiset of block sizes. We derive a
support-preserving rectangular duality, recurrences, Ehrhart evaluations, a
root-of-unity product with two degree bounds, an explicit quartic for double
chains, and a gamma-coefficient identity with both a block-boundary path
model and an algebraic constant-term derivation. For each fixed support size
$j$ the coefficient is a polynomial in $n$ of exact degree $2j$ with leading
coefficient $d^{2j}/(j!(j+1)!)$. We obtain a Stieltjes moment representation
in the chain-length variable, explicit $n^{-3/2}$ asymptotics with two
expressions for the boundary constant, and a central limit theorem for
support size including an exact variance formula and its constant-order
correction. Two independent exact verifier suites, over disjoint parameter
ranges, accompany the article.
\end{abstract}

\medskip
\noindent\textbf{Status and scope.}
The main identity is proved for every integer $d\geq1$, not just a finite
range. This is an AI-assisted research draft, not an independently refereed or
proof-assistant-verified paper; neither its correctness nor its priority has
been certified. The selected identity is stated as experimental in the
retrieved source; a targeted literature search found no subsequent proof, but
this is a bounded check and not a certification of publication priority. The
walk factorization is classical Spitzer/Wiener--Hopf background and the
earlier gamma-positivity theorem for composition polytopes is explicitly
credited; neither is claimed as new here.

\clearpage
\tableofcontents
\clearpage

\section{The selected problem and its precise meaning}
\label{sec:problem}

\subsection{Relation to the supplied manifest}
The supplied manifest~\cite{manifest} records seventy-one packages and
describes two of them on preorder polytopes: \emph{Reciprocity and Matrix
Duality for Preorder Polytopes} and \emph{A Reflexive Root-Polytope Model for
Preorder $h$-Polynomials}. Their catalogued targets concern reciprocity,
transpose duality, and polytopal realization. The present target is instead
the multiple-chain generating function in Section~9.5 of
Athanasiadis--Chapoton~\cite{AC}. This identity is not among the results
described in the manifest. No theorem claimed in those two packages is used
below; both proofs start from the defining inequalities. The manifest is a
selection and duplicate-avoidance aid, not independent evidence that its other
proofs are correct; it expressly warns that its statements describe unrefereed
claims.

\subsection{Provenance of this article}
\label{sec:provenance}
This article is the merger of two independently prepared research packages,
both dated 20~September 2026 and both targeting the same unnumbered identity.

\begin{itemize}[leftmargin=1.4em,itemsep=0.25em]
\item \texttt{enumerative-combinatorics/multiple-chain-exponential-formula}
 (the base of this merge) proved the identity by encoding each lattice point
 as a pair of subsets with a prefix-domination condition, so that the
 admissible arrays become nonnegative \emph{excursions} whose one-block symbol
 is the bounded Laurent polynomial $(1+u)^c(1+t/u)^d$; the logarithm of the
 excursion series is then the constant term $\CT P^m$, evaluated by the
 binomial theorem. It proved the statement in the rectangular generality
 $(d,c)$ recorded as Theorem~\ref{thm:main} below, and it carried the
 rectangular duality, the Ehrhart evaluation, the Stieltjes moment section,
 and the exact variance with its constant-order correction. It also contained,
 as an appendix, a rectangular form of the second proof.
\item \texttt{enumerative-combinatorics/multiple-chains-exponential-identity}
 proved the identity by running the original coordinate slack process: the
 block slack walk has step series $S_k(u)=u^k(1+q/(u-1))^k$ expanded at
 $u=\infty$, the admissible arrays are \emph{meanders} of that walk, the
 positive-part factorization is carried out in $\Q[q]((u^{-1}))[[z]]$, and
 $u=1$ is substituted only after the projection, where stars and bars
 produces the type-B Narayana polynomial. It proved the diagonal case $c=d$
 only, but it carried the fixed-support polynomiality theorem, the
 mixed-block-size theorem with its order-dependence counterexample, the
 Lagrange-inversion Narayana row, the gamma integrality and strict-positivity
 arguments, and a verifier suite over a deeper range of parameters.
\end{itemize}

\noindent
In the merged notation below, the second package's $k$ is this article's $d$
and its $q$ is this article's $t$; its $A_k(q)$ is renamed $K_d(t)$ to avoid a
collision with the array-counting symbol $\Adc n(t)$, its $\mathcal N(q,z)$ is
this article's $g(w,t)$, and its $\mathsf P_{\geq0}$ and $\mathsf P_{<0}$ are
written $[\,\cdot\,]_{\geq0}$ and $[\,\cdot\,]_{<0}$. Both packages removed
the source's Laurent normalization in the same way, and their shared numerical
data agree exactly. The second proof of Section~\ref{sec:slack} is not a
paraphrase of the first: it is a different walk, with a different symbol, in a
different ring, under a different boundary condition.

\subsection{Chain preorders, not layers of incomparable elements}
For $d\geq1$ and $n\geq0$, partition a ground set of size $dn$ into ordered
blocks $B_1,\ldots,B_n$, each of size $d$. Within a block all elements are
equivalent under the preorder, and its quotient order is
$B_1<\cdots<B_n$. The nonempty order ideals are therefore exactly the complete
prefixes $B_1\cup\cdots\cup B_r$. In particular, these are \emph{not} ordinal
sums of $d$-element antichains, whose ideals would allow partial blocks.
A vertex of the quotient poset represents an entire equivalence class, so the
preorder $\tau_{n,d}$ has $n$ quotient vertices but $dn$ underlying elements,
and its polytope has dimension $dn$, not $n$.

The associated polytope and its support enumerator are
\begin{align}
 Q_{n,d}&=\left\{a\in\R_{\geq0}^{dn}:
       \sum_{i=1}^{r}\sum_{j=1}^{d}a_{ij}\leq dr\quad(1\leq r\leq n)\right\},
       \label{eq:Q}\\
 H_{n,d}(t)&=\sum_{a\in Q_{n,d}\cap\Z^{dn}}t^{\#\{(i,j):a_{ij}>0\}},
       \qquad H_{0,d}(t)=1.\label{eq:H}
\end{align}
We use $\N=\{0,1,2,\ldots\}$. These are the source's polynomials
$h(\tau_{n,d},t)$~\cite[Sections~1, 5, and 9.5]{AC}, and~\eqref{eq:Q} is the
composition-polytope model of~\cite[Equation~(1)]{A} for the composition
$(d,d,\ldots,d)$. The support polynomial $H_{n,d}$ must not be confused with
the Ehrhart $h^*$-polynomial of $Q_{n,d}$ or with its ordinary
face-enumeration $h$-polynomial. The last inequality makes the sum finite;
the constant coefficient is $1$, and the degree is $dn$ with leading
coefficient $1$.

Put
\begin{equation}
 B_N(t)=\sum_{j=0}^N\binom Nj^2t^j,
 \qquad F_d(z,t)=\sum_{n\geq0}H_{n,d}(t)z^n.
 \label{eq:BFdefs}
\end{equation}
Specifying $B_N$ by this formula removes any ambiguity in the indexing of
Narayana polynomials: $B_N$ is the type-$\mathbb B$ Narayana polynomial in the
normalization used by the source. In a variable $q$, the source's displayed
identity reads
\begin{equation}
 1+\sum_{n\geq1}H_{n,d}(q^2)q^{-dn}x^n
 =\exp\left(\sum_{m\geq1}B_{dm}(q^2)q^{-dm}\frac{x^m}{m}\right).
 \label{eq:sourceform}
\end{equation}
The substitution $x=zq^d$ removes all negative powers. Since $t\mapsto q^2$
is injective on polynomial coefficients, the target is equivalently
\begin{equation}
 \boxed{F_d(z,t)=\exp\left(\sum_{m\geq1}B_{dm}(t)\frac{z^m}{m}\right).}
 \label{eq:target}
\end{equation}
Removing the normalization is not cosmetic: it makes the specialization $t=0$
legitimate and exposes an ordinary generating function in the number of
equivalence classes. No parity restriction on $d$ is needed.

\subsection{What is proved, and what is not claimed}
Section~\ref{sec:proof} proves~\eqref{eq:target} as a specialization of a
rectangular theorem, and Section~\ref{sec:slack} proves the same rectangular
theorem a second time by a different route.
Sections~\ref{sec:enumeration}--\ref{sec:statistics} develop its consequences.
Gamma positivity for \emph{all} composition polytopes was already proved by
Athanasiadis~\cite[Theorem~1.2]{A}; recovering it for constant block sizes is
not a new resolution of that question, and the two-subset encoding below is
closely related to that paper's encoding. The formal factorization belongs to
the classical Spitzer/Wiener--Hopf framework~\cite{Spitzer,BF}; it is not
presented as a new general lemma, and the elementary form needed here is
proved in full so that no inaccessible theorem is a hidden prerequisite. The
contribution of the proofs is to identify the exact walks to which that
factorization applies and thereby establish the selected identity.

We do not prove real-rootedness of $H_{n,d}$, gamma positivity for arbitrary
preorders, flag realizability, or the weighted $q$-reciprocity conjecture in
the same source. We do not claim that the algebraic degree bound $2^d$ is
sharp, that $H_{n,d}$ is an Ehrhart $h^*$-polynomial, that the support
polynomial of an unequal-size block sequence depends only on the multiset of
sizes (an explicit counterexample is given in
Section~\ref{sec:mixed}), or that the asymptotics of
Section~\ref{sec:asymptotics} are uniform as $d$ grows with $n$ or as
$t\to0$. Those are different assertions. In particular, a generating-function
identity with positive coefficients is not by itself a real-rootedness
argument.

\subsection{Deliberate redundancies}
\label{sec:redundancies}
Because this article merges two independent treatments, several results are
deliberately given twice. Each duplication is kept because the two versions
are not substitutes; the list below states what each one buys.

\begin{enumerate}[leftmargin=1.5em,itemsep=0.25em]
\item \textbf{Two proofs of the headline identity}
 (Sections~\ref{sec:proof} and~\ref{sec:slack}). The excursion proof gives a
 two-sided bounded symbol and a constant term evaluated by the binomial
 theorem; the slack-meander proof stays in the original coordinates, needs no
 encoding, and is the one that extends to unequal block sizes.
\item \textbf{Two forms of the factorization lemma}
 (Lemma~\ref{lem:factor} and Corollary~\ref{cor:excursion}). The Laurent-series
 form over $R((u^{-1}))$ is what the slack walk and the mixed-size theorem
 need; the constant-term excursion form is what the subset walk needs, and
 neither statement is a formal consequence of the other as stated.
\item \textbf{Two generalization directions} (Theorem~\ref{thm:main} and
 Theorem~\ref{thm:mixed}). The rectangular budget $cr$ is what powers the
 duality of Corollary~\ref{cor:duality} and the Ehrhart evaluation
 of~\eqref{eq:ehrhart}; the commutative mixed-size formula is what covers
 unequal blocks, and neither implies the other.
\item \textbf{Two derivations of the one-block Narayana series}
 (Section~\ref{sec:oneblock}). The first-step Motzkin decomposition is
 combinatorial and self-contained; the analytic route through
 $\sum_mB_m(t)w^m=\Delta^{-1/2}$ is what unlocks the Lagrange-inversion row
 \eqref{eq:d1row}.
\item \textbf{Two expressions for the asymptotic boundary constant}
 (\eqref{eq:critical} and~\eqref{eq:Kexp}). The finite radical product is the
 fast evaluation; the convergent exponential sum is what identifies the
 constant as the boundary value $F_d(\rho_d(t),t)$.
\item \textbf{Two treatments of the gamma coefficients}
 (Sections~\ref{sec:gammapath} and~\ref{sec:gammaalg}). The block-boundary
 path model supplies a combinatorial meaning, a table, and the dictionary to
 the encoding of~\cite{A}; the algebraic constant-term derivation supplies
 integrality and strict positivity, which the path model does not give.
\item \textbf{Two degree-bound arguments for $2^d$}
 (Theorem~\ref{thm:rootproduct} and Proposition~\ref{prop:degree}). The
 conjugate count is short; the explicit monic polynomial is constructive, and
 the two guard each other.
\item \textbf{Two elimination certificates for the double-chain quartic}
 (Corollary~\ref{cor:quartic} and the second pair displayed after it). They
 eliminate different auxiliary quantities and are checked by two different
 scripts.
\item \textbf{Two independent verifier suites}
 (Section~\ref{sec:checks}), over disjoint parameter ranges and with
 non-overlapping methods: only the first witnesses the subset encoding and the
 rectangular cases $c\neq d$, and only the second witnesses the endpoint
 factorization, the mixed-size recursion, and the range $d\leq8$, $n\leq12$.
\item \textbf{Two data corpora} (Appendix~\ref{app:data}). Neither is a subset
 of the other: one holds all $315$ rectangular polynomials and the only
 $\kappa_d$ table, the other reaches $d=8$ and $n=12$ and holds the only
 exact-versus-asymptotic ratio table.
\item \textbf{Two provenance records} (\texttt{sources.md}). One carries the
 hash and version reasoning for the pinned source, the other carries two
 adjacent 2026 papers and the Section~9.4 cross-check.
\end{enumerate}

One duplication has deliberately \emph{not} been kept: the weaker statement
$\Var=dn/8+O(1)$ is strictly subsumed by the exact variance of
Proposition~\ref{prop:variance} together with Theorem~\ref{thm:CLT}, so only
the sharper form is stated. Both justifications of the quasi-power step are
retained, however, in the proof of Theorem~\ref{thm:CLT}.

\section{A rectangular lattice-point theorem, proved by excursions}
\label{sec:proof}

\subsection{Statement}
For integers $d\geq1$, $c\geq0$, and $n\geq0$, let
\begin{equation}
 \mathcal A_{n;d,c}=\left\{a\in\N^{dn}:
  \sum_{i=1}^{r}\sum_{j=1}^{d}a_{ij}\leq cr\quad(1\leq r\leq n)\right\},
 \qquad
 \Adc n(t)=\sum_{a\in\mathcal A_{n;d,c}}t^{\#\supp(a)}.
 \label{eq:rectdef}
\end{equation}
The notation $\Adc0=1$ includes the empty array.

\begin{theorem}[Rectangular exponential formula]
\label{thm:main}
For every $d\geq1$ and $c\geq0$, the following identity holds in
$\Q[t][[z]]$:
\begin{equation}
 \boxed{\sum_{n\geq0}\Adc n(t)z^n
 =\exp\left(\sum_{m\geq1}\frac{z^m}{m}
       \sum_{j=0}^{\min(dm,cm)}\binom{dm}{j}\binom{cm}{j}t^j\right).}
 \label{eq:rectangular}
\end{equation}
Consequently~\eqref{eq:target} holds for every $d\geq1$. Both sides in fact
have coefficients in $\Z[t]$.
\end{theorem}

The binomial product counts a single unrestricted bridge of $m$ blocks.
The exponential is not an approximation and is not inferred from matching
initial coefficients. The next two lemmas explain exactly why it is the
series for all excursions. Section~\ref{sec:slack} proves the same theorem
again by a different walk.

\subsection{From an array to two subsets}
Flatten an array into $a_1,\ldots,a_{dn}$. Suppose its positive coordinates
occur at
\[
 S=\{s_1<\cdots<s_j\}\subseteq[dn].
\]
Their successive positive partial sums are strictly increasing:
\begin{equation}
 b_\ell=a_{s_1}+\cdots+a_{s_\ell},\qquad
 T=\{b_1<\cdots<b_j\}\subseteq[cn].
 \label{eq:ST}
\end{equation}
Here $[0]$ is the empty set. The case $j=0$ gives $S=T=\varnothing$.

\begin{lemma}[Subset encoding]
\label{lem:subsets}
The map~\eqref{eq:ST} is a support-preserving bijection between
$\mathcal A_{n;d,c}$ and pairs
\begin{equation}
 S\subseteq[dn],\quad T\subseteq[cn],\quad |S|=|T|,
 \qquad |S\cap[dr]|\leq |T\cap[cr]|\quad(1\leq r\leq n).
 \label{eq:subsetconstraint}
\end{equation}
The support size is $|S|=|T|$.
\end{lemma}
\begin{proof}
For the inverse, order $S$ and $T$ as above, put $b_0=0$, set
$a_{s_\ell}=b_\ell-b_{\ell-1}$, and put all other coordinates equal to zero.
This gives a unique nonnegative array with the specified support and partial
sums. If $k=|S\cap[dr]|$, its sum over the first $dr$ coordinates is $b_k$
when $k>0$, and zero when $k=0$. The condition $b_k\leq cr$ is equivalent to
having at least $k$ elements of $T$ in $[cr]$. This is precisely the last
inequality in~\eqref{eq:subsetconstraint}. The two constructions are inverse.
\end{proof}

Divide the positions of $S$ into $n$ blocks of length $d$, and the positions
of $T$ into $n$ blocks of length $c$. Write $p_r$ and $q_r$ for the numbers
selected in their respective $r$th blocks. The successive walk heights are
\begin{equation}
 h_r=\sum_{i=1}^r(q_i-p_i),\qquad h_0=h_n=0,\qquad h_r\geq0.
 \label{eq:heights}
\end{equation}
A block with $(p_r,q_r)=(p,q)$ has
$\binom dp\binom cq$ possible decorations, contributes weight $t^p$, and
changes the height by $q-p$. Its Laurent-polynomial symbol is therefore
\begin{equation}
 P_{d,c}(u,t)=\sum_{p=0}^d\sum_{q=0}^c
  \binom dp\binom cq t^p u^{q-p}
  =(1+u)^c(1+t/u)^d.
 \label{eq:symbol}
\end{equation}
Thus $\Adc n(t)$ is the weight of all length-$n$ nonnegative excursions
with step symbol $P_{d,c}$. Positivity is checked after a whole block;
no extra restriction is imposed on an ordering of choices inside it.

\begin{example}
Take $d=c=2$, $n=2$, and the array $(0,2\mid1,0)$. Then
$S=\{2,3\}$ and $T=\{2,3\}$. Both block heights are zero and the weight is
$t^2$. For $(0,0\mid0,4)$, the sets are $S=\{4\}$ and $T=\{4\}$; the
weight is $t$. The original coordinate partial sums and the subset-walk
heights are different statistics. Only the specified equivalence of their
prefix constraints is used.
\end{example}

\subsection{The formal factorization, in two forms}
\label{sec:factorization}
Let $R$ be any commutative algebra over $\Q$. We work with
$R((u^{-1}))$, the ring of formal Laurent series in $u$ whose exponents are
bounded above; its products are well defined, since for any fixed exponent
only finitely many pairs of exponents below the respective upper bounds
contribute. For such a series, $[\,\cdot\,]_{\geq0}$ retains its nonnegative
powers of $u$ and $[\,\cdot\,]_{<0}$ retains the strictly negative powers.
Apply these operators coefficientwise to series in $z$. At a fixed $z$-degree
only finitely many factors contribute to a product, logarithm, or exponential
with positive $z$-valuation, so all formal logarithms, exponentials,
products, and projections used below are legitimate.

The first form is stated in this full generality, with infinitely many
possible negative jumps and no positivity assumption on their weights. It is
what Section~\ref{sec:slack} and Theorem~\ref{thm:mixed} need.

\begin{lemma}[Nonnegative-walk factorization]
\label{lem:factor}
Let $S(u)\in R((u^{-1}))$. There is a unique $M(z,u)\in R[u][[z]]$ with
\begin{equation}
 M=1+z\pos{SM},\label{eq:walkrec}
\end{equation}
namely
\begin{equation}
 M(z,u)=\exp\left(\sum_{m\geq1}\frac{z^m}{m}
                                  \pos{S(u)^m}\right).
 \label{eq:meanderfactor}
\end{equation}
\end{lemma}
\begin{proof}
Set
\[
 L=-\log(1-zS)=\sum_{m\geq1}\frac{z^mS^m}{m},\qquad
 L_+=\pos L,\qquad L_-=\negp L,
\]
the projections acting coefficientwise in $z$. Each coefficient of $L_+$ is a
polynomial in $u$, so $M_*:=\exp(L_+)$ lies in $R[u][[z]]$. Since the
coefficient algebra is commutative,
\begin{equation}
 (1-zS)M_*=\exp(-L)\exp(L_+)=\exp(-L_-).\label{eq:factorization}
\end{equation}
Every monomial of $L_-$ has strictly negative $u$-exponent, hence so does
every nonconstant term of $\exp(-L_-)$; therefore
$\pos{\exp(-L_-)}=1$. Applying $\pos{\cdot}$ to~\eqref{eq:factorization},
and using that $M_*$ has only nonnegative powers of $u$, gives
$M_*-z\pos{SM_*}=1$, so $M_*$ solves~\eqref{eq:walkrec}.

Finally, writing $M=\sum_{n\geq0}M_n(u)z^n$, equation~\eqref{eq:walkrec}
determines $M_0=1$ and $M_n=\pos{SM_{n-1}}$ for $n\geq1$. This proves
existence and uniqueness, and hence~\eqref{eq:meanderfactor}.
\end{proof}

\begin{corollary}[Excursions are the exponential of bridges]
\label{cor:excursion}
Let $P(u)\in R[u,u^{-1}]$ be a Laurent \emph{polynomial} and let $E_P(z)$ be
the generating series of walks with step symbol $P$ which start and end at
height zero and remain at nonnegative heights. Then
\begin{equation}
 E_P(z)=\exp\left(\sum_{m\geq1}\frac{z^m}{m}\CT P(u)^m\right).
 \label{eq:excursionlemma}
\end{equation}
\end{corollary}
\begin{proof}
Because $P$ has a finite exponent range, each coefficient of the meander
series $M_P$ of Lemma~\ref{lem:factor} is a Laurent polynomial, and
$E_P=\CT M_P$. Now $L_+$ has only nonnegative $u$-degrees, so a product of its
monomials can have degree zero only when every chosen monomial has degree
zero; hence
\[
 \CT\exp(L_+)=\exp(\CT L_+)=\exp(\CT L),
\]
which is~\eqref{eq:excursionlemma}.
\end{proof}

\begin{remark}
\label{rem:twoforms}
Neither statement contains the other as written. Lemma~\ref{lem:factor} is
needed in $R((u^{-1}))$, where the slack symbol~\eqref{eq:slacksymbol} lives
and where the mixed-size step series of Section~\ref{sec:mixed} lives;
Corollary~\ref{cor:excursion} needs a two-sided bounded symbol for $E_P$ to be
the excursion enumerator with finitely many decorated steps per block. The
constant-term algebra in the proof of the corollary does extend verbatim to
$R((u^{-1}))$; what does not extend is the reading of $\CT P^m$ as a finite
bridge count over a bounded step set, which is exactly what the binomial
theorem evaluates in Section~\ref{sec:firstproof}.
\end{remark}

The separation of degree zero into the \emph{nonnegative} factor is essential:
the walks may touch zero at intermediate block boundaries. Replacing this by
a strictly positive convention would count a different class. One must
likewise not replace the proof by the assertion that every cyclic orbit has
exactly one nonnegative rotation: zero sums, periodic words, and unequal
positive jumps make that assertion false in this form. The factorization
handles these degeneracies without any orbit-counting assumption.

\subsection{First proof of the theorem}
\label{sec:firstproof}
\begin{proof}[First proof of Theorem~\ref{thm:main}]
Lemma~\ref{lem:subsets} and~\eqref{eq:symbol} identify the left side of
\eqref{eq:rectangular} with $E_{P_{d,c}}(z)$. The binomial theorem gives
\begin{align*}
 \CT P_{d,c}(u,t)^m
 &=\CT(1+u)^{cm}(1+t/u)^{dm}\\
 &=\sum_{j=0}^{\min(dm,cm)}\binom{cm}{j}\binom{dm}{j}t^j.
\end{align*}
Substitution in Corollary~\ref{cor:excursion}
proves~\eqref{eq:rectangular}.
For $c=d$, the defining arrays are exactly the lattice points
of~\eqref{eq:Q}, so $\Adc n=H_{n,d}$ and the selected conjecture follows.
If $c=0$, only the zero array is possible for each $n$; both sides reduce to
$(1-z)^{-1}$, consistently with the same proof. The lattice-point
interpretation places all coefficients in $\Z[t]$.
\end{proof}

\section{A second proof by the original slack process}
\label{sec:slack}

This section proves Theorem~\ref{thm:main} a second time, without the subset
encoding, by running the process on the original coordinates. The walk, the
symbol, the ring and the boundary condition are all different from those of
Section~\ref{sec:proof}: the admissible arrays are meanders with a free
terminal slack rather than excursions, the symbol is a Laurent series with an
infinite negative tail rather than a Laurent polynomial, and the logarithm is
extracted by the nonnegative projection followed by $u=1$ rather than by a
constant term. The argument also produces an endpoint refinement that the
first proof does not see.

\subsection{Slack walks and the block weight}
For an array in~\eqref{eq:rectdef}, let $Y_r=\sum_{j=1}^d a_{rj}$ and set
\[
 k_r=cr-\sum_{i=1}^rY_i,\qquad k_0=0.
\]
The inequalities say exactly that $k_r\geq0$ for every $r$; the final slack
$k_n$ may be arbitrary. The jump in slack is $c-Y_r$, bounded above by $c$,
but unbounded below before nonnegativity is imposed. Thus the arrays are
\emph{meanders}, not excursions, of this walk.

Retaining the whole contents of a block as its decoration, a block of $d$
coordinates with prescribed coordinate sum $s$ has support polynomial
\begin{equation}
 W_{d,s}(t)=
 \begin{cases}
  1,&s=0,\\
  \displaystyle\sum_{r=1}^{\min(d,s)}
       \binom{d}{r}\binom{s-1}{r-1}t^r,&s\geq1,
 \end{cases}
 \label{eq:blockweight}
\end{equation}
because the second binomial coefficient counts positive compositions of $s$
into $r$ specified coordinates. Formula~\eqref{eq:blockweight} yields a
finite-state dynamic program on the slack, using only the defining
inequalities and no conjectured identity; this is what one of the two
verifier suites of Section~\ref{sec:checks} actually runs.

\subsection{The step symbol and its ring}
The weight generating series of one coordinate is
\[
 \omega(v,t)=1+\frac{tv}{1-v}=\sum_{s\geq0}W_{1,s}(t)v^s .
\]
The slack-step symbol, expanded at $u=\infty$, is
\begin{equation}
 S_{d,c}(u,t)=u^c\,\omega(u^{-1},t)^d
             =u^c\left(\frac{u+t-1}{u-1}\right)^d
             =\sum_{s\geq0}W_{d,s}(t)\,u^{c-s}.
 \label{eq:slacksymbol}
\end{equation}
Here and below $1/(u-1)=u^{-1}+u^{-2}+u^{-3}+\cdots$ is expanded at infinity;
expansion at $u=0$ would encode a different walk and invalidate the argument.
The symbol has largest $u$-degree $c$ and an infinite negative tail, so it
lies in $\Q[t]((u^{-1}))$ and not in $\Q[t][u,u^{-1}]$. Lemma~\ref{lem:factor}
applies verbatim with $R=\Q[t]$: each $z^m$ coefficient is bounded above in
$u$-degree, and its nonnegative projection is a finite polynomial.

\subsection{The endpoint refinement}
Let
\[
 M_{d,c}(z,u,t)=\sum_{n\geq0}z^n\sum_{a\in\mathcal A_{n;d,c}}
     t^{\#\supp(a)}u^{cn-\sum_{i,j}a_{ij}}
\]
record the terminal slack as well as the support. Appending one block and
discarding negative new slack is exactly~\eqref{eq:walkrec} for
$S=S_{d,c}$, so Lemma~\ref{lem:factor} gives
\begin{equation}
 M_{d,c}(z,u,t)
 =\exp\left(\sum_{m\geq1}\frac{z^m}{m}\pos{S_{d,c}(u,t)^m}\right).
 \label{eq:slackgf}
\end{equation}

\begin{corollary}[Endpoint-refined factorization]
\label{cor:endpoint}
Since $S_{d,c}(u,t)^m=S_{dm,cm}(u,t)$, one has for $cm>0$
\begin{equation}
 \pos{S_{d,c}(u,t)^m}
 =\sum_{s=0}^{cm}W_{dm,s}(t)\,u^{cm-s}
 =u^{cm}+\sum_{j=1}^{\min(dm,cm)}\binom{dm}{j}t^j
       \sum_{\ell=j}^{cm}\binom{\ell-1}{j-1}u^{cm-\ell},
 \label{eq:slackprojection}
\end{equation}
and therefore
\begin{equation}
 M_{d,c}(z,u,t)=\exp\left(\sum_{m\geq1}\frac{z^m}{m}
       \sum_{s=0}^{cm}W_{dm,s}(t)u^{cm-s}\right).
 \label{eq:endpointfactor}
\end{equation}
When $c=0$ the projection is simply $1$.
\end{corollary}
\begin{proof}
Both displays in~\eqref{eq:slackprojection} are the truncation
of~\eqref{eq:slacksymbol} with $(d,c)$ replaced by $(dm,cm)$ to the range
$s\leq cm$; expanding $W_{dm,s}$ by~\eqref{eq:blockweight} and exchanging the
two summations gives the second form with $\ell=s$.
\end{proof}

This refinement computes the joint distribution of support and unused budget.
Its diagonal case $c=d$ is verified directly at integer $t$ by one of the
verifier suites of Section~\ref{sec:checks}.

\subsection{Evaluating the positive part}
\begin{lemma}[Unrestricted endpoint count; stars and bars]
\label{lem:starsbars}
For all integers $D\geq0$ and $C\geq0$,
\begin{equation}
 \sum_{\substack{a_1,\ldots,a_D\geq0\\ a_1+\cdots+a_D\leq C}}
   t^{\#\{i:a_i>0\}}
 =\sum_{r\geq0}\binom{D}{r}\binom{C}{r}t^r .
 \label{eq:unrestricted}
\end{equation}
\end{lemma}
\begin{proof}
Choose the $r$ positive coordinates in $\binom Dr$ ways. Subtract one from
each of them; introducing a nonnegative slack variable converts their
sum-at-most-$C$ condition into an equation with $r+1$ nonnegative unknowns and
total $C-r$, which stars and bars counts in $\binom{C}{r}$ ways. The same
formula at $r=0$ counts the zero vector. Multiplying the two choices and
summing proves the identity.
\end{proof}

\begin{proof}[Second proof of Theorem~\ref{thm:main}]
Equation~\eqref{eq:slackgf} holds by Lemma~\ref{lem:factor}. Each coefficient
of that series is a polynomial in $u$, so the substitution $u=1$ is a
well-defined ring homomorphism applied coefficientwise in $z$. The projected
polynomial at index $m$, evaluated at $u=1$, enumerates all $dm$ coordinates
with total at most $cm$; Lemma~\ref{lem:starsbars} with $D=dm$ and $C=cm$
turns this into $\sum_j\binom{dm}{j}\binom{cm}{j}t^j$. Since
$F_{d,c}(z,t)=M_{d,c}(z,1,t)$, this is~\eqref{eq:rectangular} again.
\end{proof}

Substituting $u=1$ directly into the rational
expression~\eqref{eq:slacksymbol} would be invalid, because it has a pole
there. Projecting first is what makes the proof correct. Both proofs are
independent of the manifest's reciprocity, root-polytope, and triangulation
claims: each uses only the defining prefix inequalities, a formal
factorization, and an elementary count.

\subsection{Why the type-\texorpdfstring{$\mathbb B$}{B} polynomial appears}
The coefficient in the logarithm does \emph{not} impose the intermediate
constraints at block boundaries. It imposes only nonnegative \emph{final}
slack after $m$ blocks. Those $dm$ coordinates then satisfy a single simplex
constraint, and the two binomial choices in Lemma~\ref{lem:starsbars} --
positions and values -- explain the product $\binom{dm}{j}\binom{cm}{j}$, and
hence the square when $c=d$. The exponential restores all intermediate
nonnegativity conditions through the factorization. This separation is the
structural content of the identity.

\section{Duality, recurrences, and exact enumeration}
\label{sec:enumeration}

\subsection{A support-preserving rectangular duality}
\begin{corollary}
\label{cor:duality}
For $c,d\geq1$ and all $n$, one has
\begin{equation}
 A_{n;d,c}(t)=A_{n;c,d}(t).\label{eq:duality}
\end{equation}
There is an explicit support-preserving bijection realizing this equality.
\end{corollary}
\begin{proof}
The generating-function proof follows from the symmetry of the binomial
product. For a bijection, encode an array by $(S,T)$ and put
\[
 S'=\{cn+1-b:b\in T\},\qquad
 T'=\{dn+1-s:s\in S\}.
\]
If $h_r$ is the original block height, the new height at block $r$ is
$h_{n-r}$: swap the two subset types and reverse the block sequence.
It is nonnegative, and the new path also ends at zero. Decode $(S',T')$
using Lemma~\ref{lem:subsets} with the dimensions interchanged. Applying this
map twice restores the original pair, and its common cardinality is unchanged.
\end{proof}

For example, $A_{1;2,3}(t)=1+6t+3t^2=A_{1;3,2}(t)$. The arrays live in
different ambient dimensions, so this equality is not coordinate reversal of
the original array. The subset description is what makes the bijection simple;
this corollary has no counterpart in the diagonal case alone.

\subsection{Convolution and partition formulas}
Define the bridge polynomial
\begin{equation}
 C_m^{d,c}(t)=\sum_j\binom{dm}{j}\binom{cm}{j}t^j.
 \label{eq:C}
\end{equation}
Logarithmic differentiation of~\eqref{eq:rectangular} immediately gives
\begin{equation}
 \boxed{n\Adc n(t)=\sum_{m=1}^n C_m^{d,c}(t)\Adc{n-m}(t),
        \qquad \Adc0(t)=1.}\label{eq:recurrence}
\end{equation}
Although this recursion divides by $n$, Theorem~\ref{thm:main} proves that
the division is exact coefficientwise over $\Z[t]$. The accompanying
implementations check this divisibility rather than rounding numerical values.

For a partition $\lambda\vdash n$, let $m_r(\lambda)$ be its number of parts
of size $r$ and put $z_\lambda=\prod_{r\geq1}r^{m_r}m_r!$. Expanding the
exponential also gives the finite formula
\begin{equation}
 \Adc n(t)=\sum_{\lambda\vdash n}\frac{1}{z_\lambda}
                  \prod_{r\geq1}C_r^{d,c}(t)^{m_r(\lambda)}.
 \label{eq:partition}
\end{equation}
Equivalently, because the number of permutations of cycle type $\lambda$ is
$n!/z_\lambda$,
\begin{equation}
 \Adc n(t)=\frac1{n!}\sum_{\pi\in\mathfrak S_n}
   \prod_{\gamma\text{ a cycle of }\pi}C_{|\gamma|}^{d,c}(t).
 \label{eq:permutation}
\end{equation}
This is a cycle-index identity, or weighted permutation average; it is not,
without further work, a bijection between lattice points and unlabelled
collections of decorated cycles.

In the diagonal case the first rows can be recorded symbolically in $d$:
\begin{align}
 H_{2,d}&=\frac{B_d^2+B_{2d}}{2},\qquad
 H_{3,d}=\frac{B_d^3+3B_dB_{2d}+2B_{3d}}{6},\notag\\
 H_{4,d}&=\frac{B_d^4+6B_d^2B_{2d}+3B_{2d}^2+8B_dB_{3d}+6B_{4d}}{24}.
 \label{eq:symbolicrows}
\end{align}

At $t=1$, Vandermonde's identity yields
\begin{equation}
 C_m^{d,c}(1)=\binom{(c+d)m}{dm},\qquad B_m(1)=\binom{2m}{m}.
 \label{eq:unweightedbridge}
\end{equation}
Thus the totals can be evaluated through index $N$ using $O(N^2)$ integer
arithmetic operations after precomputing the binomial coefficients. With
naive coefficient convolution, computing all support polynomials through
index $N$ takes $O(b^2N^4)$ arithmetic operations and $O(bN^2)$ stored
coefficients, where $b=\min(c,d)$ is fixed and positive. These are arithmetic
operation counts, not bit-complexity bounds.

\subsection{Examples and boundary cases}
For one block the prefix condition is just a total-sum bound, so
$H_{1,d}=B_d$. The first double-chain polynomials are
\begin{align*}
 H_{1,2}(t)&=1+4t+t^2,\\
 H_{2,2}(t)&=1+12t+27t^2+12t^3+t^4,\\
 H_{3,2}(t)&=1+24t+134t^2+236t^3+134t^4+24t^5+t^6.
\end{align*}
For instance, $2H_{2,2}=B_2^2+B_4$ follows directly from the recurrence, and
$H_{2,2}(1)=53$. These values agree with the small double-chain table
in~\cite[Section~9.4]{AC}.
Table~\ref{tab:counts} gives further totals computed independently from
coordinate dynamic programming and from~\eqref{eq:recurrence}.

\begin{table}[htbp]
\centering
\small
\begin{tabular}{r r r r r}
\toprule
$n$ & $H_{n,1}(1)$ & $H_{n,2}(1)$ & $H_{n,3}(1)$ & $H_{n,4}(1)$\\
\midrule
0&1&1&1&1\\
1&2&6&20&70\\
2&5&53&662&8885\\
3&14&554&26780&1409002\\
4&42&6362&1205961&250837850\\
5&132&77580&58050204&47928217484\\
6&429&986253&2924165436&9605146588557\\
7&1430&12927170&152231599628&1992173405162370\\
\bottomrule
\end{tabular}
\caption{Exact lattice-point counts. The $d=1$ column is the Catalan sequence
shifted by one index.}\label{tab:counts}
\end{table}

The specialization $t=0$ gives $H_{n,d}(0)=1$ and $F_d(z,0)=(1-z)^{-1}$,
as it must, since only the zero lattice point remains.
The leading coefficient of $H_{n,d}$ is also $1$: if all $dn$ coordinates are
positive, the total bound $dn$ forces every coordinate to be $1$.
These are useful normalization checks; changing the source's support
statistic or number of blocks would usually violate one of them.

\subsection{Uniform dilations and Ehrhart evaluation}
For a nonnegative integer $\ell$, the lattice points of $\ell Q_{n,d}$ are
precisely $\mathcal A_{n;d,\ell d}$. Therefore Theorem~\ref{thm:main} gives
\begin{equation}
 \sum_{n\geq0}\Ehr(Q_{n,d},\ell)z^n
 =\exp\left(\sum_{m\geq1}\binom{(\ell+1)dm}{dm}\frac{z^m}{m}\right).
 \label{eq:ehrhart}
\end{equation}
The corresponding finite partition formula is
\begin{equation}
 \Ehr(Q_{n,d},\ell)
 =\sum_{\lambda\vdash n}\frac1{z_\lambda}
       \prod_{r\geq1}\binom{(\ell+1)dr}{dr}^{m_r(\lambda)}.
 \label{eq:ehrhartpoly}
\end{equation}
For fixed $n,d$, its right side is a polynomial in $\ell$ of degree at most
$dn$. Thus this is also an identity of Ehrhart polynomials, since it holds
at every nonnegative integer $\ell$. It provides an evaluation formula in the
number of blocks; it is not a proof of real-rootedness of those polynomials,
and it does not identify the support polynomial $H_{n,d}$ with an Ehrhart
$h^*$-polynomial. This evaluation is available only because
Theorem~\ref{thm:main} is rectangular.

\section{Algebraic generating functions}
\label{sec:algebraic}

\subsection{The one-block-size series, two ways}
\label{sec:oneblock}
For $d=c=1$, the subset-walk symbol is
\begin{equation}
 P_{1,1}(u,t)=1+t+u+t/u.
 \label{eq:motzkinsymbol}
\end{equation}
These are Motzkin excursions with horizontal weight $1+t$, up weight $1$,
and down weight $t$. A nonempty excursion begins either with a horizontal
step followed by an excursion, or with an up step, an excursion translated
up one unit, a down step, and another excursion. Hence, writing
$g(w,t)=F_1(w,t)$,
\begin{equation}
 g=1+(1+t)wg+tw^2g^2.\label{eq:quadratic}
\end{equation}
The unique formal solution with constant term $1$ is
\begin{equation}
 \boxed{g(w,t)=\frac{2}{1-(1+t)w+\sqrt{\Delta(w,t)}},}
 \qquad \Delta(w,t)=1-2(1+t)w+(1-t)^2w^2.
 \label{eq:g}
\end{equation}
The square root has constant term $1$. This rationalized expression remains
valid at $t=0$, unlike a formula which divides by $t$.
At $t=1$ it becomes $g(w,1)=\Cat(w)^2$, where
$\Cat(w)=2/(1+\sqrt{1-4w})$ is the Catalan generating function, so
$H_{n,1}(1)$ is the Catalan number $\Cat_{n+1}$.

A second, independent derivation runs through the ordinary generating function
of the Narayana polynomials themselves. Summing the constant-term
expression~\eqref{eq:Bgamma} below gives
\begin{equation}
 \sum_{m\geq0}B_m(t)w^m=\frac1{\sqrt{\Delta(w,t)}},
 \label{eq:Bordinary}
\end{equation}
the square root again having constant coefficient $1$. Indeed, expand
$\bigl(1-(1+t)w-\sqrt t\,w(v+v^{-1})\bigr)^{-1}$ in powers of $v+v^{-1}$ and
take its constant term in $v$: only even powers survive, with constant term
$\binom{2r}{r}$, and summing them gives~\eqref{eq:Bordinary}. Now
Theorem~\ref{thm:main} at $d=1$ says $\log g=\sum_{m\geq1}B_m(t)w^m/m$, so
\[
 w\frac{\partial}{\partial w}\log g=\frac1{\sqrt{\Delta(w,t)}}-1,
\]
whose unique solution with constant coefficient $1$ is again~\eqref{eq:g}.

The second route is what yields a closed form for the $d=1$ row. Put
$Y=wg$; then~\eqref{eq:quadratic} becomes $Y=w(1+Y)(1+tY)$, and Lagrange
inversion gives the Narayana row
\begin{equation}
 \boxed{H_{n,1}(t)=\frac1{n+1}\sum_{j=0}^{n}
              \binom{n+1}{j}\binom{n+1}{j+1}t^j.}
 \label{eq:d1row}
\end{equation}

\subsection{A root-of-unity product for every block size}
\begin{theorem}
\label{thm:rootproduct}
Let $d\geq1$, let $\xi=e^{2\pi i/d}$, and introduce a formal variable $s$ with
$s^d=z$. Then
\begin{equation}
 \boxed{F_d(z,t)=\prod_{r=0}^{d-1}g(\xi^r s,t).}\label{eq:rootproduct}
\end{equation}
The product contains only powers of $s^d$, has rational polynomial
coefficients in $t$, and is algebraic over $\Q(z,t)$ of degree at most $2^d$.
No assertion that this upper bound is always minimal is made.
\end{theorem}
\begin{proof}
Theorem~\ref{thm:main} with $d=1$ gives
$\log g(w,t)=\sum_{m\geq1}B_m(t)w^m/m$. Therefore
\begin{align*}
 \log\prod_{r=0}^{d-1}g(\xi^rs,t)
 &=\sum_{m\geq1}\frac{B_m(t)s^m}{m}
                          \sum_{r=0}^{d-1}\xi^{rm}\\
 &=\sum_{k\geq1}\frac{B_{dk}(t)s^{dk}}{k}
 =\log F_d(s^d,t).
\end{align*}
Both series have constant coefficient $1$, so the formal exponential proves
the identity and the cancellation of fractional powers. Algebraicity follows
already from~\eqref{eq:g}, since each factor is quadratic over
$\Q(t,s)$ and algebraicity is transitive.

For the degree bound, work over the rational-function field $\Q(z,t)$ and
adjoin $s$, all $d$th roots of unity, and the $d$ square roots occurring in
the factors of~\eqref{eq:rootproduct}. This is contained in a finite normal
extension. An automorphism over $\Q(z,t)$ permutes the arguments $\xi^rs$,
and can change the signs of their square roots. As the product is insensitive
to the permutation, every conjugate is one of at most $2^d$ sign-choice
products. In characteristic zero the algebraic degree is the number of
distinct conjugates, proving the bound.
\end{proof}

\begin{proposition}[A constructive form of the same bound]
\label{prop:degree}
Let $g_r^+$ and $g_r^-$ be the two roots of the quadratic~\eqref{eq:quadratic}
at $w=\xi^rs$, in a splitting field over $\Q(z,t)$ containing all $d$ roots
$\xi^rs$ of $s^d-z$. Then
\begin{equation}
 \prod_{\varepsilon\in\{+,-\}^d}
   \left(T-\prod_{r=0}^{d-1}g_r^{\varepsilon_r}\right)
 \label{eq:monic}
\end{equation}
is a monic polynomial of degree $2^d$ with coefficients in $\Q(z,t)$, and
$F_d(z,t)$ is one of its roots.
\end{proposition}
\begin{proof}
Every automorphism of the splitting field over $\Q(z,t)$ permutes the $\xi^rs$
and, over each of them, either exchanges or preserves the two quadratic roots.
It therefore permutes the $2^d$ displayed factors, so the coefficients
of~\eqref{eq:monic} are fixed and lie in $\Q(z,t)$. The branch
in~\eqref{eq:rootproduct} is the factor with all $\varepsilon_r=+$.
\end{proof}

The conjugate count of Theorem~\ref{thm:rootproduct} is shorter; the explicit
polynomial~\eqref{eq:monic} is constructive and exhibits the annihilating
equation. Keeping both is cheap and they guard each other.

This product is a multiplicative root-of-unity filter. An \emph{average} of
the $g(\xi^rs,t)$ would instead extract the coefficients $H_{dn,1}$; those
count paths constrained after every microstep and generally give the wrong
answer. For example, $H_{1,2}(1)=6$ whereas $H_{2,1}(1)=5$.

Algebraicity also implies that $F_d$ satisfies a linear differential equation
over $\Q(z,t)$, and its coefficients in $z$ satisfy a polynomial-coefficient
recurrence in $n$. Indeed, derivatives of an algebraic function lie in the
same finite extension of $\Q(z,t)$, so sufficiently many are linearly
dependent over that field. This observation does not turn
\eqref{eq:recurrence} itself into a fixed-order, constant-coefficient recurrence.

\subsection{An explicit quartic for double chains}
\begin{corollary}
\label{cor:quartic}
The double-chain series $Y=F_2(z,t)$ is the unique solution in $\Q[t][[z]]$
with constant term $1$ of
\begin{align}
0={}&t^4z^4Y^4+t^2z^2\bigl((1+t)^2z-1\bigr)Y^3\notag\\
 &+2tz\bigl(1+(t^2+t+1)z\bigr)Y^2
    +\bigl((1+t)^2z-1\bigr)Y+1.
 \label{eq:quartic}
\end{align}
\end{corollary}
\begin{proof}
Write $z=s^2$, $a=g(s,t)$, $b=g(-s,t)$, and $Y=ab$. The two quadratic
relations~\eqref{eq:quadratic} give
\begin{align*}
 ts^2a^2+((1+t)s-1)a+1&=0,\\
 ts^2Y^2+(-(1+t)s-1)Ya+a^2&=0.
\end{align*}
Their resultant with respect to $a$, followed by $s^2=z$, is exactly the
right side of~\eqref{eq:quartic}. This polynomial identity is reproduced by
\texttt{code/derive\_quartic.py}. The defining polynomial has value $1-Y$
at $z=0$ and derivative $-1$ with respect to $Y$ there. Coefficientwise
induction, or the formal implicit-function theorem, proves uniqueness of the
branch with constant term $1$.
\end{proof}

A second elimination certificate, computed by a different script from a
different pair, reaches the same quartic. With $a=g(w,t)$ and $F=a\,g(-w,t)$,
the pair is
\begin{align}
 t w^2a^2+((1+t)w-1)a+1&=0,\notag\\
 a^2-((1+t)w+1)Fa+t w^2F^2&=0,
 \label{eq:quarticpair2}
\end{align}
whose resultant in $a$, after $w^2=z$, is again the right side
of~\eqref{eq:quartic}. The script
\texttt{code/sym\-bolic\_cer\-tif\-i\-cate.py} verifies
that this difference is the zero polynomial, rather than a series vanishing to
finite order.

At $t=1$, the equation is
\[
 z^4Y^4+(4z^3-z^2)Y^3+(6z^2+2z)Y^2+(4z-1)Y+1=0.
\]
The independent standard-library verifier also substitutes the exact support
polynomials into~\eqref{eq:quartic} through $z^{16}$, retaining every power of
$t$; this is a check of the printed certificate, not a substitute for the
resultant proof.

\section{Gamma coefficients: a path model and an algebraic derivation}
\label{sec:gamma}

The gamma expansion itself is a single identity, stated once
in~\eqref{eq:gamma} and~\eqref{eq:gammagf}. The two derivations below have
different corollaries and are both retained.

\subsection{The block-boundary path model}
\label{sec:gammapath}
For $c=d$, one has
\begin{equation}
 P_{d,d}(u,t)=(1+t+u+t/u)^d.\label{eq:groupedmotzkin}
\end{equation}
Thus an excursion with $n$ macrosteps is a word of $dn$ microsteps, grouped
into blocks of $d$. Up and down steps have increments $+1,-1$, and a
horizontal step has increment zero and two colors of weights $1,t$.
The walk must have nonnegative height only after the $d$th, $2d$th,
$\ldots$, $nd$th microsteps and must finish at zero.

\begin{proposition}
\label{prop:gamma}
Let $\gamma_{n,d,j}$ be the number of uncolored words of length $dn$ in
$\{U,D,H\}$ having $j$ up steps and $j$ down steps, total height zero,
and nonnegative heights at multiples of $d$. Then
\begin{equation}
 H_{n,d}(t)=\sum_{j=0}^{\lfloor dn/2\rfloor}
                    \gamma_{n,d,j}t^j(1+t)^{dn-2j}.
 \label{eq:gamma}
\end{equation}
In particular, $H_{n,d}$ is palindromic and unimodal.
\end{proposition}
\begin{proof}
Expand each factor of~\eqref{eq:groupedmotzkin} into its four choices.
An excursion has equally many up and down steps. If that number is $j$, the
vertical steps together contribute $t^j$, while the $dn-2j$ horizontal
positions contribute $(1+t)^{dn-2j}$. Summing over uncolored words proves
the expansion. Each displayed summand has a symmetric unimodal coefficient
sequence with center $dn/2$; their nonnegative sum has the same properties.
\end{proof}

For $d=2$, the first gamma vectors are
\begin{center}
\begin{tabular}{r l}
\toprule
$n$ & $(\gamma_{n,2,0},\gamma_{n,2,1},\ldots)$\\
\midrule
0 & $(1)$\\
1 & $(1,2)$\\
2 & $(1,8,5)$\\
3 & $(1,18,47,14)$\\
4 & $(1,32,190,244,42)$\\
5 & $(1,50,530,1560,1186,132)$\\
\bottomrule
\end{tabular}
\end{center}
For $n=1,d=2$, both $UD$ and $DU$ are allowed, explaining the gamma
coefficient $2$. The word $DU$ would be forbidden for ordinary Motzkin paths,
which require nonnegativity after every microstep.

This is a constant-block-size reformulation of the earlier composition
polytope gamma phenomenon~\cite[Theorem~1.2 and Proposition~2.2]{A}.
In the two-subset encoding, the common elements of $S$ and $T$ give the
horizontal color of weight $t$, the elements in neither give the other color,
and $T\setminus S$ and $S\setminus T$ give up and down steps. Requiring
$S\cap T=\varnothing$ leaves exactly the uncolored paths counted by
$\gamma_{n,d,j}$. This makes the connection to the prior interpretation
explicit rather than treating positivity as an uncredited new theorem.

If $\Gamma_{n,d}(y)=\sum_j\gamma_{n,d,j}y^j$, so that
\begin{equation}
 H_{n,d}(t)=(1+t)^{dn}\,\Gamma_{n,d}\!\left(\frac{t}{(1+t)^2}\right),
 \label{eq:Hgamma}
\end{equation}
the same factorization gives
\begin{equation}
 \sum_{n\geq0}\Gamma_{n,d}(y)z^n
 =\exp\left(\sum_{m\geq1}\frac{z^m}{m}G_{dm}(y)\right),
 \qquad
 G_N(y)=\sum_{r=0}^{\lfloor N/2\rfloor}\binom{N}{2r}\binom{2r}{r}y^r .
 \label{eq:gammagf}
\end{equation}
Indeed, use the uncolored symbol $(1+u+y/u)^d$; a bridge of $dm$
microsteps chooses its $2r$ vertical positions and then its $r$ up positions.

\subsection{The algebraic derivation, integrality and strict positivity}
\label{sec:gammaalg}
The same $G_N$ arises from a constant-term computation, which yields
properties the path model does not. Note first that
\begin{equation}
 \binom{N}{2r}\binom{2r}{r}=\frac{N!}{r!^2\,(N-2r)!},
\end{equation}
so that, with $v$ an auxiliary variable and $\sqrt t$ a formal symbol,
\begin{align}
 B_N(t)
 &=\CTv\bigl(1+\sqrt t\,v\bigr)^N\bigl(1+\sqrt t\,v^{-1}\bigr)^N \notag\\
 &=\sum_{r=0}^{\lfloor N/2\rfloor}\frac{N!}{r!^2(N-2r)!}\,
    t^r(1+t)^{N-2r}
 =(1+t)^NG_N\!\left(\frac{t}{(1+t)^2}\right).
 \label{eq:Bgamma}
\end{align}
The displayed constant term is a polynomial in $t$; the square root is only
an auxiliary symbol. Applying~\eqref{eq:Bgamma} factorwise
to~\eqref{eq:partition} in the diagonal case gives~\eqref{eq:Hgamma}
and~\eqref{eq:gammagf} directly, with the total factor of $1+t$ equal to
$(1+t)^{dn}$ because the parts of $\lambda$ sum to $n$.

\begin{proposition}[Integrality and strict positivity]
\label{prop:gammapos}
Every $\gamma_{n,d,j}$ is a nonnegative integer, $\gamma_{n,d,0}=1$, and for
$n\geq1$ every $\gamma_{n,d,j}$ with $0\leq j\leq\lfloor dn/2\rfloor$ is
strictly positive. Consequently the coefficients of $H_{n,d}$ increase
strictly up to the centre, except for the equal middle pair forced by odd
degree.
\end{proposition}
\begin{proof}
The exponential~\eqref{eq:gammagf} produces a priori only nonnegative rational
gamma coefficients. For integrality, observe that the coefficients of any
symmetric integer polynomial of degree $N$ in the basis $t^r(1+t)^{N-2r}$ are
integers: determine the coefficient at $r=0$, subtract that multiple, and
continue successively with the coefficient of $t^r$. Each step subtracts an
integer polynomial and requires no division. Applying this to $H_{n,d}$, whose
coefficients are integers by Theorem~\ref{thm:main}, gives integrality.

For strict positivity, expand the exponential over partitions as
in~\eqref{eq:partition}. All contributions are nonnegative, and the single-part
partition $(n)$ alone contributes $G_{dn}(y)/n$, whose coefficient of $y^j$ is
positive for every $0\leq j\leq\lfloor dn/2\rfloor$. At $y=0$ the exponential
in~\eqref{eq:gammagf} is $(1-z)^{-1}$, so $\gamma_{n,d,0}=1$; the case $n=0$ is
immediate. The final statement follows because the $r=0$ term
of~\eqref{eq:Hgamma} is $(1+t)^{dn}$, whose binomial coefficients increase
strictly to the centre, and all other terms are nonnegative and symmetric
about the same centre.
\end{proof}

Writing out the coefficient of $y$ gives the closed form
\begin{equation}
 \gamma_{n,d,1}=\frac{d^2n(n+1)}2-nd,
 \label{eq:gamma1}
\end{equation}
which also follows from the first support coefficient~\eqref{eq:h1} below by
subtracting $nd$, the coefficient of $t$ in the gamma-basis term $(1+t)^{dn}$.
As recorded in Section~\ref{sec:problem}, qualitative gamma positivity was
already proved in greater generality in~\cite{A}; what
\eqref{eq:gammagf} adds here is an explicit exponential organization in the
equal-block case, and what Proposition~\ref{prop:gammapos} adds is
integrality together with strictness.

\section{Fixed support size: polynomiality in the number of blocks}
\label{sec:fixed}

Write $h_{n,d,j}=[t^j]H_{n,d}(t)$, with $h_{n,d,j}=0$ if $j>dn$. The following
result concerns a fixed $j$, not the whole row, whose degree grows with $n$.
It has no counterpart in the rectangular development above.

\begin{theorem}[Polynomiality and leading coefficient]
\label{thm:fixed}
Fix integers $d\geq1$ and $j\geq1$. As a function of $n\in\N$, $h_{n,d,j}$ is
the restriction of a rational-coefficient polynomial of exact degree $2j$,
with leading coefficient
\begin{equation}
 \boxed{\displaystyle\frac{d^{2j}}{j!\,(j+1)!}.}
 \label{eq:leadingfixed}
\end{equation}
Equivalently, its ordinary generating function has the form
\begin{equation}
 \sum_{n\geq0}h_{n,d,j}z^n=\frac{\Pi_{d,j}(z)}{(1-z)^{2j+1}},
 \qquad\deg\Pi_{d,j}\leq2j-1,
 \label{eq:fixedrational}
\end{equation}
with
\begin{equation}
 \Pi_{d,j}(1)=\Cat_j\,d^{2j},\qquad
 \Cat_j=\frac1{j+1}\binom{2j}{j}.
 \label{eq:Pboundary}
\end{equation}
In particular the pole order $2j+1$ is exact.
\end{theorem}
\begin{proof}
Separate the constant-in-$t$ part of the logarithm in~\eqref{eq:target}:
\begin{equation}
 F_d(z,t)=\frac1{1-z}
   \exp\left(\sum_{r\geq1}t^r\Lambda_{d,r}(z)\right),
 \qquad
 \Lambda_{d,r}(z)=\sum_{m\geq1}\binom{dm}{r}^{\!2}\frac{z^m}{m}.
 \label{eq:fixedstart}
\end{equation}
For $r\geq1$, the expression $\binom{dm}{r}^2/m$ is a polynomial in $m$ of
degree $2r-1$: indeed $\binom{dm}{r}$ has a factor $m$, so the quotient is
polynomial and has zero constant term. The elementary identity
\[
 \sum_{m\geq1}m^az^m=\left(z\frac{\mathrm d}{\mathrm dz}\right)^a\frac{z}{1-z}
\]
shows that $\Lambda_{d,r}$ has denominator dividing $(1-z)^{2r}$ and a
numerator of degree at most $2r-1$.

In the coefficient of $t^j$ in the exponential, a product
$\prod_r\Lambda_{d,r}^{e_r}$ satisfies $\sum_r re_r=j$; its denominator
divides $(1-z)^{2j}$ and its numerator has degree at most
$2j-\sum_re_r\leq2j-1$. The extra factor $(1-z)^{-1}$
in~\eqref{eq:fixedstart} proves~\eqref{eq:fixedrational}.

For the highest pole, the leading term of $\binom{dm}{r}^2/m$ is
$d^{2r}m^{2r-1}/(r!)^2$, so
\begin{equation}
 \Lambda_{d,r}(z)=\frac{d^{2r}(2r-1)!}{(r!)^2}(1-z)^{-2r}
   +O\bigl((1-z)^{-2r+1}\bigr).
 \label{eq:Aleading}
\end{equation}
The coefficient of the highest pole in~\eqref{eq:fixedrational} is therefore
$[\sigma^j]$ of
\[
 \exp\left(\sum_{r\geq1}\frac1{2r}\binom{2r}{r}(d^2\sigma)^r\right)
 =\Cat(d^2\sigma),
\]
the last identity following, for instance, by differentiating $\log\Cat(v)$
and using $\Cat(v)=2/(1+\sqrt{1-4v})$. This proves~\eqref{eq:Pboundary}.

Finally, each monomial $z^a$ in $\Pi_{d,j}$ contributes
$\binom{n+2j-a}{2j}$ to the coefficient of $z^n$. Since $a\leq2j-1$, this
binomial expression is a polynomial in $n$ agreeing with the coefficient for
every $n\geq0$, including the initial zero values, and its leading
coefficient is $1/(2j)!$. Hence the leading coefficient of $h_{n,d,j}$ is
$\Pi_{d,j}(1)/(2j)!=d^{2j}/(j!(j+1)!)$.
\end{proof}

\subsection{The first two coefficients}
At support size one, either~\eqref{eq:fixedstart} or a direct coordinate count
gives
\begin{equation}
 h_{n,d,1}=d^2\binom{n+1}{2}.
 \label{eq:h1}
\end{equation}
For the direct argument, a coordinate in block $b$ may contain any positive
integer at most $bd$ if all other coordinates vanish; there are $d$ positions
in that block, and $\sum_{b=1}^nd(bd)=d^2\binom{n+1}{2}$.

At support size two,
\begin{align}
 h_{n,d,2}
 &=\frac{d^2(d-1)^2}{4}\binom{n+3}{4}
   +\frac{3d^4-d^2}{2}\binom{n+2}{4}\notag\\
 &\hspace{3em}+\frac{d^2(d+1)^2}{4}\binom{n+1}{4}.
 \label{eq:h2}
\end{align}
To check this directly, use
$[t^2]F_d=(1-z)^{-1}\bigl(\Lambda_{d,2}+\tfrac12\Lambda_{d,1}^2\bigr)$; put it
over $(1-z)^5$, whose numerator is
\[
 \tfrac14\bigl(d^4(z+6z^2+z^3)-2d^3(z-z^3)+d^2(z-2z^2+z^3)\bigr),
\]
and extract coefficients. The leading coefficients of~\eqref{eq:h1}
and~\eqref{eq:h2} are $d^2/2$ and $d^4/12$, in agreement with
Theorem~\ref{thm:fixed}.

\section{A Stieltjes moment representation in the number of blocks}
\label{sec:moments}

Fix a real number $t>0$. The variable in this section is the number of
blocks $n$, not the support degree. Let $U$ be the bilateral shift on
$\ell^2(\Z)$ and define
\begin{equation}
 J=(1+t)I+\sqrt t\,(U+U^*)
    =(I+\sqrt t\,U)(I+\sqrt t\,U^*).
 \label{eq:J}
\end{equation}
This is a bounded positive self-adjoint operator. If $P$ is the orthogonal
projection onto $\ell^2(\N)$, let
\begin{equation}
 T_d=PJ^dP\big|_{\ell^2(\N)}.\label{eq:compression}
\end{equation}
The compression is applied \emph{after} the $d$th power. In general,
$PJ^dP$ and $(PJP)^d$ are different: the latter would disallow negative
intermediate microstep heights.

\begin{theorem}
\label{thm:moment}
For every $d\geq1$ and $t>0$,
\begin{equation}
 H_{n,d}(t)=\langle e_0,T_d^n e_0\rangle\qquad(n\geq0).
 \label{eq:moment}
\end{equation}
Consequently there is a probability measure $\mu_{d,t}$ supported on
$[0,(1+\sqrt t)^{2d}]$ such that
\begin{equation}
 H_{n,d}(t)=\int x^n\,d\mu_{d,t}(x).
 \label{eq:measure}
\end{equation}
Every leading Hankel determinant
$\det(H_{i+j,d}(t))_{0\leq i,j\leq r}$ is strictly positive. In particular,
$H_{n,d}(t)^2<H_{n-1,d}(t)H_{n+1,d}(t)$ for all $n\geq1$.
\end{theorem}
\begin{proof}
Matrix multiplication expands~\eqref{eq:moment} as a sum over the
block-boundary paths of Section~\ref{sec:gammapath}, with each vertical
microstep of weight $\sqrt t$ and each horizontal one of weight $1+t$.
For a returning path with $j$ up and $j$ down steps, the vertical weight is
$t^j$, the same as before. This proves~\eqref{eq:moment}.

The Fourier multiplier of $J$ is
$1+t+2\sqrt t\cos\theta$, which lies between $0$ and $(1+\sqrt t)^2$.
Thus $0\leq T_d\leq(1+\sqrt t)^{2d}I$. The spectral theorem applied to
$T_d$ and the unit vector $e_0$ gives~\eqref{eq:measure} with the stated
support and total mass.

The Hankel matrix is the Gram matrix of
$e_0,T_de_0,\ldots,T_d^re_0$. The vector $T_d^i e_0$ is supported on
$\{0,\ldots,di\}$ and has coefficient $t^{di/2}>0$ at $di$.
These vectors are therefore linearly independent by comparing their last
nonzero coordinates. Their Gram determinant is strictly positive.
In particular, the representing measure has infinite support. Applying strict
Cauchy--Schwarz to $x^{(n-1)/2}$ and $x^{(n+1)/2}$ in this measure proves
the strict log-convexity statement.
\end{proof}

The shifted Hankel matrices are strictly positive definite as well.
For a nonzero finitely supported $v$, Fourier transformation shows
\[
 \langle v,T_dv\rangle
 =\frac1{2\pi}\int_0^{2\pi}
 (1+t+2\sqrt t\cos\theta)^d|\widehat v(\theta)|^2\,d\theta>0.
\]
The multiplier is positive except possibly at a single point when $t=1$.
Applying this to linear combinations of $T_d^ie_0$ proves the assertion.
At $t=0$, by contrast, $H_{n,d}(0)=1$ and the Hankel matrices have rank one.
No coefficientwise Hankel-positivity assertion in the indeterminate $t$ is
needed or claimed here.

\section{Unequal block sizes}
\label{sec:mixed}

The factorization of Lemma~\ref{lem:factor} does not require every step to
have the same block size. There is, however, an essential distinction between
a \emph{fixed order} of unequal sizes and a \emph{sum over all orders}. This
section states the positive result and the counterexample which shows that its
qualification cannot be dropped. Throughout it, the budget equals the block
size, so this generalizes the diagonal case $c=d$ only; the rectangular
direction of Theorem~\ref{thm:main} and the mixed-size direction here are
orthogonal, and neither implies the other.

Let $\mathcal D$ be a finite set of positive integers. For a word
$\kappa=(d_1,\ldots,d_r)$ over $\mathcal D$, define $H_\kappa(t)$ by the
prefix constraints at the partial sums $d_1+\cdots+d_i$, exactly as
in~\eqref{eq:Q}; the empty word has polynomial $1$. Introduce commuting
variables $x_d$ for $d\in\mathcal D$ and set
\[
 \mathcal F_{\mathcal D}(t;\boldsymbol x)
 =\sum_{\kappa\text{ a word over }\mathcal D}
    H_\kappa(t)\prod_{i=1}^{|\kappa|}x_{d_i}.
\]
For a nonnegative multiplicity vector $\boldsymbol m=(m_d)_{d\in\mathcal D}$,
put $|\boldsymbol m|=\sum_dm_d$ and $\|\boldsymbol m\|=\sum_ddm_d$.

\begin{theorem}[Commutative mixed-size formula]
\label{thm:mixed}
One has
\begin{equation}
 \boxed{\mathcal F_{\mathcal D}(t;\boldsymbol x)
 =\exp\left(\sum_{\boldsymbol m\neq0}
   \frac{(|\boldsymbol m|-1)!}{\prod_{d\in\mathcal D}m_d!}
   B_{\|\boldsymbol m\|}(t)\prod_{d\in\mathcal D}x_d^{m_d}\right).}
 \label{eq:mixed}
\end{equation}
For fixed multiplicities, the left side sums the polynomials of all
\emph{distinct ordered words} with those multiplicities.
\end{theorem}
\begin{proof}
Replace $zS_d$ in Lemma~\ref{lem:factor} by
$\mathcal S(u)=\sum_{d\in\mathcal D}x_dS_{d,d}(u,t)$; all formal operations
are justified by positive total degree in the $x_d$. The same argument gives
the exponential of $\pos{\sum_{r\geq1}\mathcal S^r/r}$. In a monomial with
multiplicities $\boldsymbol m$, the multinomial coefficient is
$|\boldsymbol m|!/\prod_dm_d!$, and by~\eqref{eq:slacksymbol}
\[
 \prod_{d\in\mathcal D}S_{d,d}(u,t)^{m_d}
 =S_{\|\boldsymbol m\|,\|\boldsymbol m\|}(u,t).
\]
Project to nonnegative powers, set $u=1$, and apply
Lemma~\ref{lem:starsbars} with $D=C=\|\boldsymbol m\|$. Dividing the
multinomial coefficient by $|\boldsymbol m|$ gives~\eqref{eq:mixed}.
\end{proof}

It would be incorrect to infer that $H_\kappa$ depends only on the multiset of
its parts. Direct enumeration gives
\begin{align}
 H_{(1,1,2)}(t)&=1+11t+23t^2+11t^3+t^4,\notag\\
 H_{(1,2,1)}(t)&=1+11t+22t^2+11t^3+t^4,\label{eq:ordercounterexample}\\
 H_{(2,1,1)}(t)&=1+11t+23t^2+11t^3+t^4.\notag
\end{align}
The mixed formula gives their sum as the coefficient of $x_1^2x_2$. This
example also explains why the constant-block problem is special: a multiset of
identical parts admits only one order. It is what makes the phrase ``summed
over orders'' in Theorem~\ref{thm:mixed} a genuine restriction rather than a
formality.

\begin{remark}[The reach of the scalar method]
\label{rem:reach}
A fixed periodic pattern of unequal sizes may be grouped into a macrostep only
if its internal block-boundary restrictions are retained; simply replacing the
period by its total size erases inequalities and is generally invalid, as
\eqref{eq:ordercounterexample} shows. It does \emph{not} follow that the
method stops at constant block sizes. Theorem~\ref{thm:mixed} shows that the
same scalar factorization covers every finite set of block sizes, with the
one qualification that its exponential computes the sum over all distinct
orders with given multiplicities rather than any individual $H_\kappa$. What
remains genuinely out of reach of the present argument is a closed formula for
a single prescribed order of unequal sizes.
\end{remark}

\section{Singularity analysis and explicit leading asymptotics}
\label{sec:asymptotics}

Fix $d\geq1$ and $t>0$. Define
\begin{equation}
 R(t)=(1+\sqrt t)^2,\qquad \rho_d(t)=R(t)^{-d},\qquad
 K_d(t)=\prod_{r=0}^{d-1}g\!\left(\frac{\xi^r}{R(t)},t\right).
 \label{eq:critical}
\end{equation}
The factors are obtained by continuation from their series at zero. In the
factor with $r=0$, the square root in~\eqref{eq:g} vanishes and the value is
$R(t)/\sqrt t$. Conjugate factors pair, and $K_d(t)$ is a positive real
number. It is also the finite boundary value $F_d(\rho_d(t),t)$, and admits a
second, convergent expression
\begin{equation}
 K_d(t)=\exp\left(\sum_{m\geq1}\frac{B_{dm}(t)}{m\,R(t)^{dm}}\right).
 \label{eq:Kexp}
\end{equation}
Both are used below: the finite radical product of~\eqref{eq:critical} is the
fast evaluation, while the exponential sum~\eqref{eq:Kexp} is what identifies
the constant as the boundary value of $F_d$. Convergence of~\eqref{eq:Kexp}
follows from the coefficient asymptotics of the Narayana polynomial itself,
\begin{equation}
 B_N(t)=\frac{1+\sqrt t}{2\sqrt\pi\,t^{1/4}}\,R(t)^NN^{-1/2}
        \bigl(1+O(N^{-1})\bigr),
 \label{eq:Basymp}
\end{equation}
which is the square-root transfer applied to~\eqref{eq:Bordinary} and makes
the summands of~\eqref{eq:Kexp} of order $m^{-3/2}$.

\begin{theorem}[Uniform-chain asymptotics]
\label{thm:asymptotics}
For fixed $d\geq1$ and $t>0$,
\begin{equation}
 \boxed{H_{n,d}(t)=
 \frac{K_d(t)(1+\sqrt t)}{2\sqrt{\pi d}\,t^{1/4}}
 (1+\sqrt t)^{2dn}n^{-3/2}\left(1+O(n^{-1})\right).}
 \label{eq:asymptotic}
\end{equation}
The radius of convergence of $F_d(\cdot,t)$ is $\rho_d(t)$, and
$K_d(t)=F_d(\rho_d(t),t)$. In particular,
\begin{equation}
 H_{n,d}(1)=\frac{K_d(1)}{\sqrt{\pi d}}\,
                4^{dn}n^{-3/2}\left(1+O(n^{-1})\right),
 \qquad
 K_d(1)=\prod_{r=0}^{d-1}
                 \left(\frac{2}{1+\sqrt{1-\xi^r}}\right)^2.
 \label{eq:totalasymptotic}
\end{equation}
The square roots in the last formula are principal, with the $r=0$ value
zero. The error in~\eqref{eq:asymptotic} is locally uniform in a complex
neighborhood of any fixed positive $t$, using the corresponding analytic
branches; in particular it is locally uniform near $t=1$.
\end{theorem}
\begin{proof}
The discriminant factors as
\begin{equation}
 \Delta(w,t)=(1-R(t)w)\bigl(1-(1-\sqrt t)^2w\bigr).
 \label{eq:deltafact}
\end{equation}
Its first positive root is $w_0=1/R(t)$; the other root is strictly farther
from zero, or is absent when $t=1$. There are no additional nonzero poles
of~\eqref{eq:g}: if its denominator vanished, multiplying it by the
conjugate expression would give $4tw^2=0$, while at $w=0$ the denominator
is $2$.

Write $w=w_0(1-\varepsilon)$. Since
\[
 1-\frac{(1-\sqrt t)^2}{(1+\sqrt t)^2}
 =\frac{4\sqrt t}{(1+\sqrt t)^2},
\]
expanding the square root and denominator in~\eqref{eq:g} gives
\begin{equation}
 g(w,t)=g(w_0,t)
 \left(1-\frac{1+\sqrt t}{t^{1/4}}\sqrt{\varepsilon}
       +O(\varepsilon)\right),
 \qquad g(w_0,t)=\frac{R(t)}{\sqrt t}.
 \label{eq:baseexpansion}
\end{equation}
Choose the local branch $s=z^{1/d}$ with $s=w_0$ at $z=\rho_d(t)$.
Only the factor $g(s,t)$ in~\eqref{eq:rootproduct} is singular there;
all $g(\xi^rs,t)$ with $r\ne0$ are analytic and nonzero. Moreover,
\[
 1-s/w_0=\frac1d(1-z/\rho_d(t))+O((1-z/\rho_d(t))^2).
\]
It follows that
\begin{equation}
 F_d(z,t)=K_d(t)
 \left(1-\frac{1+\sqrt t}{t^{1/4}\sqrt d}
               \sqrt{1-z/\rho_d(t)}+O(1-z/\rho_d(t))\right).
 \label{eq:localF}
\end{equation}
In fact the algebraic expression has a full convergent local Puiseux
expansion, not just the displayed two terms. Its leading value is
$K_d(t)>0$, so the square-root term does not cancel, and it agrees
with~\eqref{eq:Kexp} by the convergent exponential formula.

To check dominance, a singularity of a factor at
$\xi^rs=w_0$ always projects to the same point $z=w_0^d$.
The second discriminant root projects outside this circle. Rotation
invariance removes the apparent branching at $z=0$.
Thus there is a unique singularity on $|z|=\rho_d(t)$, with analytic
continuation to a dented neighborhood beyond that circle. The same separation
persists under sufficiently small complex perturbations of $t$.

Finally,
\[
 [z^n](1-z/\rho)^{1/2}
 =-\frac{\rho^{-n}}{2\sqrt\pi\,n^{3/2}}
                         \left(1+O(n^{-1})\right).
\]
The next nonanalytic Puiseux term has exponent $3/2$ and contributes
$O(\rho^{-n}n^{-5/2})$; the locally analytic terms are handled by the
continuation just established. This is the usual algebraic singularity
transfer underlying directed-path asymptotics~\cite{BF,FO}.
Applying it to~\eqref{eq:localF} proves~\eqref{eq:asymptotic}, including the
uniform form. Specialization $t=1$ gives~\eqref{eq:totalasymptotic}.
\end{proof}

At $d=1$, $K_1(1)=4$, recovering the asymptotics of the shifted Catalan
numbers. At $d=2$,
\begin{equation}
 K_2(1)=\frac{16}{(1+\sqrt2)^2}=16(3-2\sqrt2)=48-32\sqrt2,
 \quad
 H_{n,2}(1)\sim\frac{16(3-2\sqrt2)}{\sqrt{2\pi}}\,
                    16^n n^{-3/2}.
 \label{eq:doubleconstant}
\end{equation}
The two displayed closed forms are the radical product and the value obtained
from~\eqref{eq:Kexp}; they are the same number.
The exponent $-3/2$ is characteristic of excursions; despite its exponential
presentation, the generating function is algebraic rather than entire.

\begin{table}[htbp]
\centering
\begin{tabular}{r r r r}
\toprule
$d$ & $K_d(1)$ & $K_d(1)/\sqrt{\pi d}$ & growth factor $4^d$\\
\midrule
1&4.0000000000&2.2567583342&4\\
2&2.7451660041&1.0951627857&16\\
3&2.3004650571&0.7493420036&64\\
4&2.0653267193&0.5826179108&256\\
5&1.9170238358&0.4836905186&1024\\
\bottomrule
\end{tabular}
\caption{Decimal approximations to explicit algebraic boundary values and
leading asymptotic constants. Numerical values are not interval-certified.}
\label{tab:constants}
\end{table}

\begin{table}[htbp]
\centering
\begin{tabular}{r r r r}
\toprule
$d$ & $K_d(1)$ & ratio at $n=100$ & ratio at $n=400$\\
\midrule
1 & 4.0000000000 & 0.9743200738 & 0.9934736873\\
2 & 2.7451660041 & 0.9848655014 & 0.9961813732\\
3 & 2.3004650571 & 0.9886892314 & 0.9971536342\\
4 & 2.0653267193 & 0.9907300253 & 0.9976704816\\
8 & 1.6768915769 & 0.9941248072 & 0.9985270518\\
\bottomrule
\end{tabular}
\caption{Boundary constants and ratios of exact totals to the leading term
in~\eqref{eq:totalasymptotic}. Ratios approaching $1$ illustrate, but do not
prove, the asymptotic theorem.}
\label{tab:ratios}
\end{table}

The asymptotic statement does not include $t=0$. At that boundary
$H_{n,d}(0)=1$ and the square-root regime degenerates to a simple pole. Nor is
any uniformity claimed when $d$ grows with $n$, or when $t$ tends to zero with
$n$.

\section{Support size: exact variance and a central limit theorem}
\label{sec:statistics}

Choose a lattice point uniformly from $Q_{n,d}\cap\Z^{dn}$ and let $X_{n,d}$
be its number of nonzero coordinates. Its probability generating function is
$H_{n,d}(t)/H_{n,d}(1)$. Write $a_{n,d}=H_{n,d}(1)$.

\begin{proposition}[Exact mean and variance]
\label{prop:variance}
For every $d\geq1$ and $n\geq0$,
\begin{equation}
 \E X_{n,d}=\frac{dn}{2},\qquad
 \boxed{\Var X_{n,d}=\frac{d}{2a_{n,d}}
       \sum_{m=1}^n\binom{2dm-2}{dm-1}a_{n-m,d}.}
 \label{eq:exactvariance}
\end{equation}
The empty sum gives variance zero when $n=0$.
\end{proposition}
\begin{proof}
Palindromicity, that is $H_{n,d}(t)=t^{dn}H_{n,d}(t^{-1})$, makes the
distribution symmetric about $dn/2$ and gives the exact mean. For the
variance, introduce the centered series
\[
 G(z,s)=\sum_{n\geq0}e^{-dns/2}H_{n,d}(e^s)z^n
       =F_d(ze^{-ds/2},e^s).
\]
By~\eqref{eq:target},
\[
 \log G(z,s)=\sum_{m\geq1}\frac{z^m}{m}
                 e^{-dms/2}B_{dm}(e^s).
\]
Every centered polynomial on the right has first derivative zero at $s=0$.
For $N\geq1$, Vandermonde's identity and its first two marked versions give
\begin{equation}
 \sum_{j=0}^N\binom Nj^2(j-N/2)^2
 =\binom{2N}{N}\frac{N^2}{4(2N-1)}.
 \label{eq:bridgevariance}
\end{equation}
For completeness, $B_N'(1)=N\binom{2N-1}{N-1}$ and
$B_N''(1)=N(N-1)\binom{2N-2}{N-2}$; substituting these into
$\sum j^2\binom Nj^2-N\sum j\binom Nj^2+(N^2/4)B_N(1)$ proves
\eqref{eq:bridgevariance}, with $N=1$ also immediate.

Since $\partial_s\log G(z,0)=0$, one has
$\partial_s^2G(z,0)=G(z,0)\partial_s^2\log G(z,0)$.
When $N=dm$, the factor in~\eqref{eq:bridgevariance} divided by $m$ simplifies
to
\[
 \frac1m\binom{2dm}{dm}\frac{(dm)^2}{4(2dm-1)}
       =\frac d2\binom{2dm-2}{dm-1}.
\]
Taking the coefficient of $z^n$ gives~\eqref{eq:exactvariance}.
\end{proof}

\begin{theorem}[Gaussian support fluctuations]
\label{thm:CLT}
For fixed $d\geq1$, as $n\to\infty$,
\begin{equation}
 \frac{X_{n,d}-dn/2}{\sqrt{dn/8}}\ \Longrightarrow\ \mathcal N(0,1),
 \qquad
 \Var X_{n,d}=\frac{dn}{8}+\kappa_d+O(n^{-1}),
 \label{eq:CLT}
\end{equation}
where
\begin{equation}
 \boxed{\kappa_d=\frac1{16}+\frac d8
                -\frac18\sum_{r=0}^{d-1}\sqrt{1-\xi^r}.}
 \label{eq:kappa}
\end{equation}
The root sum is real; conjugate terms pair and the $r=0$ term is zero.
For example, $\kappa_1=3/16$ and
$\kappa_2=5/16-\sqrt2/8$.
\end{theorem}
\begin{proof}
Let $C_d(t)$ denote the leading prefactor in~\eqref{eq:asymptotic}.
Its locally uniform analytic form implies, for complex $s$ near zero,
\begin{equation}
 \log\E e^{sX_{n,d}}
 =n\lambda_d(s)+\log\frac{C_d(e^s)}{C_d(1)}+O(n^{-1}),
 \qquad
 \lambda_d(s)=2d\log\frac{1+e^{s/2}}2.
 \label{eq:quasipower}
\end{equation}
Two independent justifications of this quasi-power step are available, and
each suffices. The first is the locally uniform complex-$t$ error statement of
Theorem~\ref{thm:asymptotics}. The second tracks the distinguished branch
point explicitly: choosing $\sqrt t$ analytic near $t=1$, the branch point
$1/R(t)$ varies analytically while the second root of~\eqref{eq:deltafact}
stays outside a larger disk, and the $d-1$ nondistinguished factors
of~\eqref{eq:rootproduct} remain analytic and nonzero on a common
neighborhood of their evaluation points; hence the local Puiseux coefficients
and the transfer estimates vary analytically and uniformly, and
$C_d(t)=K_d(t)(1+\sqrt t)/(2\sqrt{\pi d}\,t^{1/4})$ is analytic and nonzero
there. Either way the error and its derivatives are controlled on a smaller
neighborhood by Cauchy's formula.

The derivatives of the rate at zero are $\lambda_d'(0)=d/2$ and
$\lambda_d''(0)=d/8$. Substitution $s=iv/\sqrt{dn/8}$, followed by subtraction
of the exact mean, makes~\eqref{eq:quasipower} converge to $-v^2/2$;
convergence of characteristic functions proves the central limit theorem.
Twice differentiating gives the variance expansion, with
$\kappa_d=(\log C_d(e^s))''|_{s=0}$.

It remains to evaluate this derivative. The non-$K_d$ part of the prefactor is
\[
 \frac{1+e^{s/2}}{2e^{s/4}\sqrt{\pi d}}
 =\frac{\cosh(s/4)}{\sqrt{\pi d}},
\]
which contributes $1/16$. Put $v=\tanh(s/4)$ and $q=\xi^r$. At the critical
argument $w=q/R(e^s)$, formula~\eqref{eq:g} becomes
\begin{equation}
 g(w,e^s)=
 \frac{2}{1-\tfrac12(1+v^2)q+\sqrt{1-q}\sqrt{1-v^2q}}.
 \label{eq:criticalfactor}
\end{equation}
At $v=0$ the denominator is
$(1+\sqrt{1-q})^2/2$. Its expansion in $v^2$ shows that the coefficient
of $v^2$ in $\log g(w,e^s)$ is
\[
 \frac{q}{1+\sqrt{1-q}}=1-\sqrt{1-q}.
\]
Since $v^2=s^2/16+O(s^4)$, the second derivative is
$(1-\sqrt{1-q})/8$. Summing over $r$ and adding $1/16$ proves
\eqref{eq:kappa}.
\end{proof}

An equivalent positive, convergent series is
\begin{equation}
 \kappa_d=\frac1{16}+\frac d8
       \sum_{m\geq1}\frac{\binom{2dm}{dm}}{4^{dm}(2dm-1)}.
 \label{eq:kappaseries}
\end{equation}
To see this, expand
$\sqrt{1-z}=1-\sum_{k\geq1}\binom{2k}{k}z^k/(4^k(2k-1))$
and sum at the $d$th roots of unity. The series is absolutely convergent even
at $z=1$, so the filtering is justified. In particular, $\kappa_d>1/16$.

For the double chain, the exact centered variance at $n=400$ is
approximately $0.13552348$ after subtracting $dn/8=100$; the limiting
correction is $0.13572330$. The error is consistent with the proved
$O(n^{-1})$ expansion. This numerical comparison is illustrative, whereas
the constant is determined by the finite radical expression above.

This argument concerns the uniform distribution on \emph{lattice points}, not
uniform sampling of faces, of equivalence classes, or of support sets. Those
would be different probability models.

\section{Reproducibility, proof audit, and remaining questions}
\label{sec:verification}

\subsection{Independent checks actually performed}
\label{sec:checks}
The package contains a single standard-library Python verifier which runs
\emph{two separate suites}, inherited from the two merged packages. They use
different methods over disjoint parameter ranges, and neither range has been
narrowed. All checks use arbitrary-precision integers or exact rational
numbers, and the verifier raises explicit exceptions rather than relying on
\texttt{assert}, so its checks remain active under \texttt{python -O}.

\medskip
\noindent\textbf{Suite~1 (rectangular).}
Three structurally different methods compute the same polynomials. The first
performs dynamic programming on the original coordinates, merging states with
equal current total and support size; it never uses the bridge formula. The
second uses the decorated subset-walk macrosteps of~\eqref{eq:symbol}, and is
the only computational witness of the encoding of Lemma~\ref{lem:subsets}. The
third uses the exponential-formula recurrence~\eqref{eq:recurrence} with
exact divisibility checks. All three agree for the $315$ parameter triples
\[
 1\leq d\leq5,\qquad0\leq c\leq6,\qquad0\leq n\leq8.
\]
The same run visits $34\,689$ feasible arrays individually in $48$ smaller
cases, without merged states, and matches their support distributions.
It verifies $225$ instances of rectangular duality, $54$ independent gamma
path expansions, and $54$ exact variance identities against
\eqref{eq:exactvariance}. It also checks the quartic through $z^{16}$ and
matches the first four double-chain polynomials printed in the source.

\medskip
\noindent\textbf{Suite~2 (diagonal, deeper).}
The second suite computes $H_{n,d}$ from the block-weight slack dynamic
program of~\eqref{eq:blockweight}, which uses the defining inequalities
directly, and compares it with the recurrence over
$1\leq d\leq8$, $0\leq n\leq12$: $104$ complete polynomials, well beyond the
range of Suite~1. It re-enumerates every individual coordinate vector in the
$28$ cases with $dn\leq8$; verifies the endpoint-refined
factorization~\eqref{eq:endpointfactor} in $112$ cases with
$1\leq d\leq4$, $0\leq n\leq6$ and $t\in\{0,1,2,7\}$ -- the only computational
witness of Corollary~\ref{cor:endpoint}; reconstructs the gamma basis on all
$104$ rows and checks strict positivity; confirms
Theorem~\ref{thm:fixed} by finite differences in $40$ cases
($1\leq d\leq8$, $1\leq j\leq5$), checking both the exact degree $2j$ and the
leading coefficient~\eqref{eq:leadingfixed}; checks~\eqref{eq:h1}
and~\eqref{eq:h2} on all $104$ rows; re-checks the quartic through $z^{16}$
coefficientwise in $t$; and verifies the mixed-size
identity~\eqref{eq:mixed} over all $121$ ordered compositions with parts in
$\{1,2,3\}$ and at most four blocks, against $35$ multiplicity vectors,
recording the counterexample~\eqref{eq:ordercounterexample}.

\medskip
The optional SymPy scripts derive the quartic by two different resultants,
independently of the recurrence; the second reports that its difference from
the printed quartic is the zero polynomial rather than a truncated series.
The optional mpmath script computes exact integer totals, and exact rational
means and variances, before any decimal conversion, then compares them with
the asymptotic formulas at $80$-digit working precision. It records
observations up to $n=400$ for $d\in\{1,2,3,4,5,8\}$, the exact-to-asymptotic
ratios of Table~\ref{tab:ratios}, the $\kappa_d$ comparisons, and samples at
the positive support weight $t=4$. Those decimal calculations are not
certified interval computations.

Finite checks corroborate the proofs; they do not certify untested parameters.

\subsection{A short audit of the mathematical dependencies}
Both proofs of the selected identity use only formal commutative series
algebra together with an elementary count: the first uses the subset bijection
and the binomial theorem, the second uses the slack process and stars and
bars. Their only infinite objects are formal power series and Laurent series
with exponents bounded above, handled coefficientwise. No appeal is made to
convergence or numerical extrapolation. The source's normalization is removed
explicitly in~\eqref{eq:sourceform}. The main formal results are independent
of both the analytic consequences and the computer calculations.

The algebraic formulas follow from a quadratic equation and a finite
root-of-unity filter. The degree bound is an upper bound, proved twice, and
the quartic has its formal branch specified by a nonzero derivative. The
moment representation uses the spectral theorem for a bounded positive
operator. The asymptotic and limit-law sections additionally use analytic
continuation, square-root singularity transfer~\cite{FO}, and the continuity
theorem for characteristic functions; their hypotheses are checked in place.
These analytic conclusions are separate from, and are not needed to prove, the
original conjectural identity.

Particular pitfalls which are avoided are: imposing bounds inside a block;
identifying a preorder-equivalence block with an antichain; confusing support
$h$ with Ehrhart $h^*$; counting strict rather than weak excursions; using an
additive instead of multiplicative root-of-unity filter; replacing the
compressed power $PJ^dP$ by the power of a compression; substituting $u=1$
into the slack symbol before projecting; and assuming that a cyclic orbit has
exactly one nonnegative rotation.

\subsection{Reproduction commands}
From the package directory, the complete exact check is
\begin{lstlisting}
python code/verify.py
\end{lstlisting}
It needs only Python 3.10 or later and no third-party package. By default it
writes to \texttt{results/}; \texttt{--output DIRECTORY} redirects it. All
generated coefficient arrays use increasing powers. The optional symbolic and
numerical checks are
\begin{lstlisting}
python code/derive_quartic.py
python code/symbolic_certificate.py
python code/asymptotics.py
\end{lstlisting}
The first two require SymPy and the third mpmath; \texttt{asymptotics.py}
additionally accepts \texttt{--maximum N} (default $400$). Versions used in
this run are recorded in \texttt{results/environment.json}. The article can be
rebuilt with
\begin{lstlisting}
latexmk -pdf -interaction=nonstopmode article.tex
\end{lstlisting}
or with the two \texttt{pdflatex} passes performed by \texttt{build.sh} and by
the optional \texttt{Makefile}.
The original supplied manifest is preserved in \texttt{context/manifest.tex};
the two merged packages shipped byte-identical copies of it, and only one is
kept. The proposed additional entry in \texttt{manifest-entry.tex} describes
this package without modifying that source file.

\subsection{What would be a genuinely different next problem?}
The all-block-size identity selected here no longer has an unresolved case
within either proof, and Theorem~\ref{thm:mixed} extends the method to
arbitrary finite sets of block sizes once one sums over orders. What remains
open is of a different kind. A closed formula for a single prescribed order of
unequal sizes is not delivered by the present scalar factorization, and
\eqref{eq:ordercounterexample} shows that no formula depending only on the
multiset can exist. Real-rootedness of the support polynomials, the minimal
algebraic degree for a given $d$ (only the upper bound $2^d$ is proved here),
gamma positivity for arbitrary preorders, and the weighted $q$-reciprocity
conjecture of the same source all require additional arguments. Finally, a
direct bijection explaining the permutation average~\eqref{eq:permutation},
with all periodicity factors treated explicitly, would be a useful complement
to the formal proofs.

\appendix
\section{A compact proof certificate}
\label{app:certificate}

The logical core of the second proof can be recorded in six equations, so that
the main result can be audited without reading any consequence section. All
projections act in $\Q[t]((u^{-1}))[[z]]$:
\begin{align*}
 S_{d,c}(u,t)&=u^c\left(\frac{u+t-1}{u-1}\right)^d,
       & M&=1+z\pos{S_{d,c}M},\\
 L&=-\log(1-zS_{d,c})=L_++L_-,
       & M_*&=\exp(L_+),\\
 (1-zS_{d,c})M_*&=\exp(-L_-),
       &\pos{\exp(-L_-)}&=1.
\end{align*}
The recurrence uniquely gives $M=M_*$. Since
\[
 \left.\pos{S_{d,c}^m}\right|_{u=1}
 =\sum_{a_1+\cdots+a_{dm}\leq cm}t^{\#\{i:a_i>0\}}
 =\sum_{r\geq0}\binom{dm}{r}\binom{cm}{r}t^r,
\]
specializing $u=1$ proves~\eqref{eq:rectangular}, and $c=d$ gives the
conjectured formula. The only infinite series in $u$ have exponents bounded
above, and the specialization is made only after projection, coefficient by
coefficient in $z$.

\section{Exact data and source provenance}
\label{app:data}

\begin{longtable}{@{}>{\raggedright\arraybackslash}p{0.44\textwidth}
                    >{\raggedright\arraybackslash}p{0.52\textwidth}@{}}
\toprule
File & Contents\\
\midrule
\endhead
\texttt{results/verification.json} & Exact test ranges, counts of cases, and
pass status for both verifier suites.\\
\texttt{results/exact\_polynomials.json} & All $315$ rectangular support
polynomials in increasing coefficient order; the only data with $c\neq d$.\\
\texttt{results/coefficients.json} & Every diagonal polynomial for
$1\leq d\leq8$, $0\leq n\leq12$, with its gamma coefficients; the only data
reaching $d=8$ and $n=12$.\\
\texttt{results/gamma\_polynomials.json} & Gamma coefficients obtained by
independent uncolored microstep dynamic programming.\\
\texttt{results/small\_values.csv} & Uniform-chain totals, support
coefficients, and exact rational variances.\\
\texttt{results/lattice\_point\_totals.csv} & Totals at $t=1$ over the deeper
diagonal range.\\
\texttt{results/quartic\_certificate.txt} & The first symbolic resultant, the
printed quartic, and the branch check.\\
\texttt{results/symbolic\_certificate.json} & The second elimination
certificate, with a zero-polynomial residual.\\
\texttt{results/asymptotics.csv} & Exact-count versus asymptotic ratios and
variance-correction comparisons.\\
\texttt{results/asymptotic\_constants.csv},
\texttt{results/asymptotic\_checks.csv} & Boundary constants and
exact-to-leading-term ratios at $n=100$ and $n=400$, including $d=8$.\\
\texttt{results/constants.json} & Decimal values of $K_d(1)$, leading
constants, and $\kappa_d$ for $1\leq d\leq5$; the $\kappa_d$ column occurs
nowhere else.\\
\texttt{results/pdf\_validation.json} & Machine-readable render audit of the
compiled article.\\
\texttt{results/environment.json} & Exact Python, SymPy, mpmath, and pdfTeX
versions.\\
\texttt{sources.md}, \texttt{STATUS.md} & Source/version audit and explicit
claim boundaries.\\
\bottomrule
\end{longtable}

The target was checked against the text of Section~9.5, printed page~24 of
the arXiv v1 PDF, and the corresponding HTML formula. The mathematical formula
was read in structured HTML and parsed PDF text; repeated attempts to obtain a
screenshot of that page failed with retrieval and cache errors, and no
unavailable screenshot is treated as evidence.
The source paper is not included in this archive; retrieval details are
recorded in \texttt{sources.md}. The supplied manifest is included unaltered,
with its hash recorded in \texttt{sources.md}.

\begin{thebibliography}{9}
\raggedright

\bibitem{AC}
C.~A. Athanasiadis and F. Chapoton,
\emph{Polytopes and posets associated to preorders},
arXiv:2605.26916v1 (2026), especially Section~9.5, p.~24.
\href{https://arxiv.org/abs/2605.26916v1}{arXiv record};
\href{https://arxiv.org/html/2605.26916v1}{HTML text}.
Accessed 20 September 2026.

\bibitem{A}
C.~A. Athanasiadis,
\emph{Lattice point enumeration of polytopes associated to integer compositions},
arXiv:2510.23903v1 (2025), Equation~(1), Theorem~1.2 and Section~2.
\href{https://arxiv.org/abs/2510.23903v1}{arXiv record}.
Accessed 20 September 2026.

\bibitem{Spitzer}
F. Spitzer,
A combinatorial lemma and its application to probability theory,
\emph{Transactions of the American Mathematical Society}
\textbf{82} (1956), 323--339.
\href{https://doi.org/10.1090/S0002-9947-1956-0079851-X}
{doi:10.1090/S0002-9947-1956-0079851-X}.

\bibitem{BF}
C. Banderier and P. Flajolet,
Basic analytic combinatorics of directed lattice paths,
\emph{Theoretical Computer Science} \textbf{281} (2002), 37--80.
\href{https://doi.org/10.1016/S0304-3975(02)00007-5}
{doi:10.1016/S0304-3975(02)00007-5}.

\bibitem{FO}
P. Flajolet and A. Odlyzko,
Singularity analysis of generating functions,
\emph{SIAM Journal on Discrete Mathematics} \textbf{3} (1990), no.~2,
216--240.
\href{https://doi.org/10.1137/0403019}{doi:10.1137/0403019}.

\bibitem{DHLTW}
Z. Dai, Q. Hou, Z. Liu, W. Thawinrak, and H. Wang,
\emph{Counting Lattice Points in Minkowski Sums of Cross Polytopes},
arXiv:2608.16037v2 (2026).
\href{https://arxiv.org/html/2608.16037v2}{HTML text}.
Consulted as adjacent preorder-support and duality literature; its
introduction connects support duality to Conjecture~5.4 of~\cite{AC}.

\bibitem{Yan}
X. Yan,
\emph{Variations of colored multiset Eulerian polynomials and applications},
arXiv:2608.15682v2 (2026).
\href{https://arxiv.org/html/2608.15682v2}{HTML text}.
Consulted as adjacent composition-polytope literature.

\bibitem{manifest}
V. Reshetnikov,
\emph{manifest.tex}, user-supplied \LaTeX{} research catalogue of
seventy-one packages, received 20 September 2026. The catalogue records
claimed results, not independent proof verification. A copy of the exact
supplied file is included in \texttt{context/manifest.tex}.

\end{thebibliography}
\end{document}
